\section{Θέματα Β}
\def\xrwma{secondcolor}
%=========== 1ο ΘΕΜΑ Β ===========
\begin{thema}{Β}
Δίνεται η συνάρτηση $f(x)=\dfrac{x}{e^x}$.
\begin{erwthma}
\item Να μελετήσετε την $f$ ως προς τη μονοτονία και τα ακρότατα και να βρεθεί το σύνολο τιμών της.
\item Μελετήστε την $f$ ως προς την κυρτότητα και τα σημεία καμπής.
\item Να βρεθεί η εξίσωση της εφαπτομένης της $C_f$ στο σημείο $A(0,f(0))$.
\item Να βρεθεί το εμβαδόν του χωρίου που περικλείεται από την $C_f$ τον άξονα $x'x$ τον άξονα $y'y$ και την ευθεία $x=\ln{9}$.
\end{erwthma}
\end{thema}
%=========== ΤΕΛΟΣ - 1ο ΘΕΜΑ Β ===========
%=========== 2ο ΘΕΜΑ Β ===========
\begin{thema}{Β}
Δίνονται οι συναρτήσεις $f(x)=\ln{\left(1+\frac{1}{x}\right)}$ και $g(x)=e^x$.
\begin{erwthma}
\item Να ορίσετε τη συνάρτηση $h=f\circ g$.
\end{erwthma}
Αν $h(x)=\ln\left(1+\frac{1}{e^x}\right)$ με $x\in\mathbb{R}$:
\begin{erwthma}[resume]
\item δείξτε ότι η $h$ αντιστρέφεται και ότι $h^{-1}(x)=-\ln{\left(e^x-1\right)}\ ,\ x>0$.
\item βρείτε τις ασύμπτωτες των γραφικών παραστάσεων των συναρτήσεων $h$ και $h^{-1}$.
\item βρείτε την εξίσωση της εφαπτομένης της $C_h$ η οποία είναι παράλληλη με την ευθεία $\zeta: y=-\frac{x}{2}+1$.
\end{erwthma}
\end{thema}
%=========== ΤΕΛΟΣ - 2ο ΘΕΜΑ Β ===========
%=========== 3ο ΘΕΜΑ Β ===========
\begin{thema}{Β}
Δίνεται η συνάρτηση $f(x)=\dfrac{2x}{x^2+1}$.
\begin{erwthma}
\item Να μελετήσετε την $f$ ως προς τη μονοτονία και τα ακρότατα.
\item Μελετήστε την $f$ ως προς την κυρτότητα και τα σημεία καμπής.
\item Να βρεθούν οι ασύμπτωτες της $C_f$ και να χαράξετε τη γραφική της παράσταση.
\item Να βρεθεί το εμβαδόν του χωρίου που περικλείεται από την $C_f$ τον άξονα $x'x$ και τις ευθείες $x=-1,x=2$.
\end{erwthma}
\end{thema}
%=========== ΤΕΛΟΣ - 3ο ΘΕΜΑ Β ===========
%=========== 4ο ΘΕΜΑ Β ===========
\begin{thema}{Β}
Δίνεται η συνάρτηση $f:(0,+\infty)\to\mathbb{R}$ με τύπο $f(x)=\sqrt{x}-\ln{x}$.
\begin{erwthma}
\item Να μελετήσετε την $f$ ως προς τη μονοτονία και τα ακρότατα.
\item Μελετήστε την $f$ ως προς την κυρτότητα.
\item Να βρεθεί η εξίσωση της εφαπτομένης $\varepsilon$ της $C_f$ στο σημείο $A(1,f(1))$.
\item Αν $\varepsilon: y=-\frac{x}{2}+\frac{3}{2}$ είναι η εξίσωση της εφαπτομένης στο $A$, να δείξετε ότι ισχύει $2\sqrt{x}+x-3\geq 2\ln{x}$ για κάθε $x\in(0,16]$.
\item Υπολογίστε το όριο $\lim\limits_{x\to0^+}{(xf(x))}$.
\item Να βρεθεί το εμβαδόν του χωρίου που περικλείεται από την $C_f$ την ευθεία $\varepsilon$ και την ευθεία $x=3$.
\end{erwthma}
\end{thema}
%=========== ΤΕΛΟΣ - 4ο ΘΕΜΑ Β ===========

%=========== 5ο ΘΕΜΑ Β ===========
\begin{thema}{Β}
Δίνονται οι συναρτήσεις $f:(0,+\infty)\to\mathbb{R}$ με $f(x)=x\ln{x}-x$ και $g:\mathbb{R}\to\mathbb{R}$ με $g(x)=e^x$.
\begin{erwthma}
\item Να ορίσετε τη συνάρτηση $h=f\circ g$.
\end{erwthma}
Αν $h(x)=e^x(x-1)\ ,\ x\in\mathbb{R}$:
\begin{erwthma}[resume]
\item μελετήστε την $h$ ως προς τη μονοτονία και τα ακρότατα.
\item να μελετήσετε την $h$ ως προς την κυρτότητα και τα σημεία καμπής.
\item να βρεθεί το εμβαδόν του χωρίου που περικλείεται από την $C_h$, τον άξονα $x'x$ και την ευθεία $x=-2$.
\end{erwthma}
\end{thema}
%=========== ΤΕΛΟΣ - 5ο ΘΕΜΑ Β ===========
%=========== 6ο ΘΕΜΑ Β ===========
\begin{thema}{Β}
Δίνονται οι συναρτήσεις $g:(1,+\infty)\to\mathbb{R}$ με τύπο $g(x)=\ln{\left(x-\frac{1}{x}\right)}$ και $h:\mathbb{R}\to\mathbb{R}$ με $h(x)=e^x$.
\begin{erwthma}
\item Να αποδείξετε ότι η συνάρτηση $f=h\circ g$ έχει πεδίο ορισμού $D_f=(1,+\infty)$ και τύπο $f(x)=x-\frac{1}{x}$.
\item Εξετάστε αν οι συναρτήσεις $f$ και $k$ με $k(x)=x-\frac{1}{x}$ είναι ίσες. Σε περίπτωση που δεν είναι ίσες, να βρείτε το ευρύτερο υποσύνολο του $\mathbb{R}$ στο οποίο ισχύει $f(x)=k(x)$.
\item Δείξτε ότι η $f$ είναι αντιστρέψιμη και βρείτε την $f^{-1}$.
\item Αν $f^{-1}(x)=\frac{x+\sqrt{x^2+4}}{2}\ ,\ x>0$ να υπολογίσετε το όριο
\[ \lim\limits_{x\to+\infty}{\frac{f^{-1}(x)}{x}} \]
\end{erwthma}
\end{thema}
%=========== ΤΕΛΟΣ - 6ο ΘΕΜΑ Β ===========

%=========== 7ο ΘΕΜΑ Β ===========
\begin{thema}{Β}
Δίνεται η συνάρτηση $f: \mathbb{R} \to \mathbb{R}$ με τύπο $f(x) = (x-a)e^x + 1$, όπου $a \in \mathbb{R}$. Δίνεται ακόμα ότι η συνάρτηση $f$ παρουσιάζει τοπικό ακρότατο στο σημείο $x_0 = 1$.

\begin{erwthma}
\item Να αποδείξετε ότι $a = 2$. 

Στα ερωτήματα Β2 έως Β4, να θεωρήσετε ότι $a = 2$.
\item Να μελετήσετε τη συνάρτηση $f$ ως προς τη μονοτονία και να αποδείξετε ότι η εξίσωση $f(x) = 0$ έχει ακριβώς δύο πραγματικές ρίζες. 
\item Να μελετήσετε τη συνάρτηση $f$ ως προς την κυρτότητα, να βρείτε το σημείο καμπής της γραφικής της παράστασης και να αποδείξετε ότι η ευθεία $y = 1$ είναι οριζόντια ασύμπτωτη της $C_f$ στο $-\infty$. 
\item Να υπολογίσετε το εμβαδόν του χωρίου που περικλείεται από τη γραφική παράσταση της συνάρτησης $f$, τον άξονα $x'x$ και τις κατακόρυφες ευθείες $x = 2$ και $x = 3$. 
\end{erwthma}
\end{thema}
%=========== ΤΕΛΟΣ - 7ο ΘΕΜΑ Β ===========
%=========== 8ο ΘΕΜΑ Β ===========
\begin{thema}{Β}
Δίνεται η συνάρτηση $f:(0, +\infty) \to \mathbb{R}$ με τύπο $f(x) = \dfrac{\ln x}{x}$.
\begin{erwthma}
\item Να μελετήσετε τη συνάρτηση $f$ ως προς τη μονοτονία και να βρείτε τα τοπικά ακρότατα.
\item Να μελετήσετε τη συνάρτηση $f$ ως προς την κυρτότητα και να βρείτε τις συντεταγμένες του σημείου καμπής της γραφικής της παράστασης.
\item Να βρείτε τις ασύμπτωτες της γραφικής παράστασης της συνάρτησης $f$.
\item Να υπολογίσετε το εμβαδόν του χωρίου που περικλείεται από τη γραφική παράσταση της συνάρτησης $f$, τον άξονα $x'x$ και τις κατακόρυφες ευθείες $x = 1$ και $x = e$.
\end{erwthma}
\end{thema}
%=========== ΤΕΛΟΣ - 8ο ΘΕΜΑ Β ===========

%=========== 9ο ΘΕΜΑ Β ===========
\begin{thema}{Β}
Δίνονται οι συναρτήσεις $f(x) = \frac{x}{x+1}$ με πεδίο ορισμού $D_f = \mathbb{R} - \{-1\}$ και $g(x) = e^x$ με πεδίο ορισμού $A_g = \mathbb{R}$.
\begin{erwthma}
\item Να αποδείξετε ότι η συνάρτηση $h = f \circ g$ έχει πεδίο ορισμού το $\mathbb{R}$ και τύπο $h(x) = \frac{e^x}{e^x+1}$.
\item Να αποδείξετε ότι η συνάρτηση $h$ αντιστρέφεται και να βρείτε την αντίστροφή της, $h^{-1}$.
\item Να υπολογίσετε το όριο: $\lim\limits_{x \to 1^-} \frac{h^{-1}(x)}{\ln(1-x)}$.
\item Να αποδείξετε ότι η εξίσωση $h(x) = x$ έχει τουλάχιστον μία ρίζα στο διάστημα $(0, 1)$.
\end{erwthma}
\end{thema}
%=========== ΤΕΛΟΣ - 9ο ΘΕΜΑ Β ===========
%=========== 10ο ΘΕΜΑ Β ===========
\begin{thema}{Β}
Δίνεται η συνάρτηση $f(x) = \dfrac{\ln x + 2}{x}$.
\begin{erwthma}
\item Να βρείτε το πεδίο ορισμού της συνάρτησης $f$ και να υπολογίσετε τα όρια $\lim\limits_{x \to 0^+} f(x)$ και $\lim\limits_{x \to +\infty} f(x)$. Στη συνέχεια, να προσδιορίσετε τις ασύμπτωτες της γραφικής παράστασης της $f$.
\item Να μελετήσετε τη συνάρτηση $f$ ως προς τη μονοτονία και να βρείτε το ολικό ακρότατό της.
\item Να μελετήσετε τη συνάρτηση $f$ ως προς την κυρτότητα και να βρείτε το σημείο καμπής της $C_f$.
\item Να υπολογίσετε το εμβαδόν του χωρίου που περικλείεται από τη γραφική παράσταση της συνάρτησης $f$, τον άξονα $x'x$ και τις κατακόρυφες ευθείες $x = 1$ και $x = e$.
\end{erwthma}
\end{thema}
%=========== ΤΕΛΟΣ - 10ο ΘΕΜΑ Β ===========
%=========== 11ο ΘΕΜΑ Β ===========
\begin{thema}{Β}
Δίνονται οι συναρτήσεις $f(x) = e^x + x - 1$ με πεδίο ορισμού το $\mathbb{R}$ και $g(x) = \ln x$ με πεδίο ορισμού το $(0, +\infty)$.
\begin{erwthma}
\item Να βρείτε το πεδίο ορισμού και τον τύπο της συνάρτησης $h = f \circ g$.
\item Να υπολογίσετε το όριο: $\lim\limits_{x \to 1^+} \frac{2x+1}{h(x)}$.
\item Να αποδείξετε ότι υπάρχει τουλάχιστον ένας αριθμός $x_0 \in (1, e)$ τέτοιος ώστε να ισχύει $x_0 + \ln x_0 = 2$.
\item Να αποδείξετε ότι $\displaystyle\int_{1}^{2} h(x) dx < 1 + \ln 2$.
\end{erwthma}
\end{thema}
%=========== ΤΕΛΟΣ - 11ο ΘΕΜΑ Β ===========
%=========== 12ο ΘΕΜΑ Β ===========
\begin{thema}{Β}
Δίνεται η συνάρτηση $f(x) = e^{x-1} + x^2$ με πεδίο ορισμού το $\mathbb{R}$.
\begin{erwthma}
\item Να μελετήσετε τη συνάρτηση $f$ ως προς την κυρτότητα.
\item Να βρείτε την εξίσωση της εφαπτομένης της γραφικής παράστασης της συνάρτησης $f$ στο σημείο της με τετμημένη $x_0 = 1$.
\item Να αποδείξετε ότι $e^{x-1} -3x \ge -x^2 - 1$ για κάθε $x \in \mathbb{R}$ και να βρείτε πότε ισχύει η ισότητα.
\item Να αποδείξετε ότι $\displaystyle\int_{0}^{2} f(x) dx > 4$.
\end{erwthma}
\end{thema}
%=========== ΤΕΛΟΣ - 12ο ΘΕΜΑ Β ===========

%=========== 13ο ΘΕΜΑ Β ===========
\begin{thema}{Β}

Δίνεται η συνάρτηση $f(x) = \dfrac{a x + 1}{x - 2}$ με πεδίο ορισμού το $D_f = \mathbb{R} - \{2\}$, της οποίας η γραφική παράσταση διέρχεται από το σημείο $M(3, 4)$. 
\begin{erwthma}
\item Να αποδείξετε ότι $a = 1$.
\item Να βρείτε τις κατακόρυφες και οριζόντιες ασύμπτωτες της γραφικής παράστασης της συνάρτησης $f$.
\item Αν $g(x) = \ln x$, να βρείτε το πεδίο ορισμού και τον τύπο της συνάρτησης $h = f \circ g$.
\item Να αποδείξετε ότι η εξίσωση $h(x) = 2026$ έχει τουλάχιστον μία ρίζα στο διάστημα $(e^2, e^3)$.
\end{erwthma}
\end{thema}
%=========== ΤΕΛΟΣ - 13ο ΘΕΜΑ Β ===========

%=========== 14ο ΘΕΜΑ Β ===========
\begin{thema}{Β}
Δίνονται οι συναρτήσεις $f(x) = x^3 + x - 2$ και $g(x) = e^{x-1}$, ορισμένες στο $\mathbb{R}$.
\begin{erwthma}
\item Να βρείτε τον τύπο της συνάρτησης $h = f \circ g$.
\item Να αποδείξετε ότι η συνάρτηση $h$ είναι «1-1» και να λύσετε την εξίσωση $h(x^2) = h(2x-1)$.
\item Να βρείτε την εξίσωση της εφαπτομένης της γραφικής παράστασης της συνάρτησης $h$ η οποία είναι παράλληλη στην ευθεία $\varepsilon: y = 4x + 2026$.
\item Να υπολογίσετε το εμβαδόν του χωρίου που περικλείεται από τη γραφική παράσταση της συνάρτησης $f$, τον άξονα $x'x$ και τις κατακόρυφες ευθείες $x=0$ και $x=1$.
\end{erwthma}
\end{thema}
%=========== ΤΕΛΟΣ - 14ο ΘΕΜΑ Β ===========


%=========== 15ο ΘΕΜΑ Β ===========
\begin{thema}{Β}

Δίνεται η συνάρτηση $f(x) = \ln(x - a) + x - 1$. Δίνεται επίσης ότι η γραφική παράσταση της $f$ διέρχεται από το σημείο $A(2, 1)$.
\begin{erwthma}
\item Να αποδείξετε ότι $a = 1$ και να βρείτε το πεδίο ορισμού της $f$.
\item Να βρείτε την κατακόρυφη ασύμπτωτη της γραφικής παράστασης της $f$.
\item Να αποδείξετε ότι η $f$ αντιστρέφεται και να λύσετε την εξίσωση $f(e^x + 1) = 1$.
\item Να αποδείξετε ότι η εξίσωση $f(x) = 0$ έχει ακριβώς μία ρίζα στο διάστημα $(1, 2)$.
\end{erwthma}
\end{thema}
%=========== ΤΕΛΟΣ - 15ο ΘΕΜΑ Β ===========
%=========== 16ο ΘΕΜΑ Β ===========
\begin{thema}{Β}
Δίνεται η συνάρτηση $f(x) = \dfrac{x^2+1}{x}$ με $x \ne 0$.
\begin{erwthma}
\item Να μελετήσετε την $f$ ως προς τη μονοτονία και τα ακρότατα.
\item Να μελετήσετε την $f$ ως προς την κυρτότητα και να δείξετε ότι η γραφική της παράσταση δεν έχει σημεία καμπής.
\item Να βρείτε τις ασύμπτωτες της γραφικής παράστασης της συνάρτησης $f$.
\item Να υπολογίσετε το εμβαδόν του χωρίου που περικλείεται από τη γραφική παράσταση της $f$, την πλάγια ασύμπτωτη της $C_f$ στο $+\infty$ και τις ευθείες $x=1$ και $x=e$.
\end{erwthma}
\end{thema}
%=========== ΤΕΛΟΣ - 16ο ΘΕΜΑ Β ===========

%=========== 17ο ΘΕΜΑ Β ===========
\begin{thema}{Β}
Δίνεται η συνάρτηση $f(x) = \dfrac{e^x}{x^2}$ με $x \ne 0$.
\begin{erwthma}
\item Να βρείτε τα όρια $\lim\limits_{x \to 0} f(x)$ και $\lim\limits_{x \to -\infty} f(x)$. Στη συνέχεια να βρείτε, αν υπάρχουν, τις οριζόντιες και κατακόρυφες ασύμπτωτες της $C_f$.
\item Να μελετήσετε τη συνάρτηση $f$ ως προς τη μονοτονία και να βρείτε τα τοπικά ακρότατα.
\item Να βρείτε την εξίσωση της εφαπτομένης της $C_f$ στο σημείο της $A(-1, f(-1))$.
\item Να αποδείξετε ότι ο ισχυρισμός «υπάρχει μοναδικό $x_0 \in (1, 2)$ τέτοιο ώστε $f(x_0)=2$» είναι ψευδής.
\end{erwthma}
\end{thema}
%=========== ΤΕΛΟΣ - 17ο ΘΕΜΑ Β ===========

%=========== 18ο ΘΕΜΑ Β ===========
\begin{thema}{Β}
Δίνεται η συνάρτηση $f(x) = (x-2)\ln(x-2) + x - 2$.
\begin{erwthma}
\item Να βρείτε το πεδίο ορισμού της $f$ και το όριο $\lim\limits_{x \to 2^+} f(x)$.
\item Να μελετήσετε την $f$ ως προς τη μονοτονία και τα τοπικά ακρότατα.
\item Να μελετήσετε την $f$ ως προς την κυρτότητα.
\item Να υπολογίσετε το εμβαδόν του χωρίου που περικλείεται από τη γραφική παράσταση της $f$, τον άξονα $x'x$ και τις ευθείες $x=3$ και $x=e+2$.
\end{erwthma}
\end{thema}
%=========== ΤΕΛΟΣ - 18ο ΘΕΜΑ Β ===========

%=========== 19ο ΘΕΜΑ Β ===========
\begin{thema}{Β}
Δίνεται η συνάρτηση $f(x) = a x^3 - 3x + \beta$, όπου $a, \beta \in \mathbb{R}$. Δίνεται ότι η γραφική παράσταση της $f$ διέρχεται από το σημείο $A(0, 1)$ και η $f$ παρουσιάζει τοπικό ακρότατο στο $x_0 = 1$.
\begin{erwthma}
\item Να αποδείξετε ότι $a=1$ και $\beta=1$.
\item Να μελετήσετε την $f$ ως προς τη μονοτονία και την κυρτότητα.
\item Να βρείτε την εξίσωση της εφαπτομένης της $C_f$ στο σημείο της καμπής της.
\item Να αποδείξετε ότι η εξίσωση $f(x) = 0$ έχει ακριβώς μία πραγματική ρίζα, η οποία ανήκει στο διάστημα $(-2, -1)$.
\end{erwthma}
\end{thema}
%=========== ΤΕΛΟΣ - 19ο ΘΕΜΑ Β ===========

%=========== 20ο ΘΕΜΑ Β ===========
\begin{thema}{Β}
Δίνεται η συνάρτηση $f(x) = \dfrac{e^x}{e^x-1}$.
\begin{erwthma}
\item Να βρείτε το πεδίο ορισμού της $f$ καθώς και τα όρια $\lim\limits_{x \to 0^+} f(x)$ και $\lim\limits_{x \to +\infty} f(x)$.
\item Να αποδείξετε ότι η συνάρτηση $f$ αντιστρέφεται και να βρείτε το πεδίο ορισμού και τον τύπο της αντίστροφης συνάρτησης $f^{-1}$.
\item Να βρείτε τις ασύμπτωτες της γραφικής παράστασης της $f^{-1}$.
\item Να υπολογίσετε το εμβαδόν του χωρίου που περικλείεται από τη γραφική παράσταση της $f$ (για $x>0$), τον άξονα $x'x$ και τις ευθείες $x=\ln 2$ και $x=\ln 3$.
\end{erwthma}
\end{thema}
%=========== ΤΕΛΟΣ - 20ο ΘΕΜΑ Β ===========

%=========== 21ο ΘΕΜΑ Β ===========
\begin{thema}{Β}
Δίνεται η συνάρτηση $f(x) = x^2 e^{-x}$.
\begin{erwthma}
\item Να μελετήσετε τη συνάρτηση $f$ ως προς τη μονοτονία και να βρείτε τα ακρότατά της.
\item Να μελετήσετε τη συνάρτηση $f$ ως προς την κυρτότητα και να βρείτε τα σημεία καμπής της $C_f$.
\item Να βρείτε τις ασύμπτωτες της γραφικής παράστασης της $f$.
\item Να υπολογίσετε το εμβαδόν του χωρίου που περικλείεται από τη γραφική παράσταση της $f$, τον άξονα $x'x$ και τις ευθείες $x=0$ και $x=1$.
\end{erwthma}
\end{thema}
%=========== ΤΕΛΟΣ - 21ο ΘΕΜΑ Β ===========

%=========== 22ο ΘΕΜΑ Β ===========
\begin{thema}{Β}
Δίνεται η συνάρτηση $f(x) = \ln(x+1) - \dfrac{x}{x+1}$.
\begin{erwthma}
\item Να βρείτε το πεδίο ορισμού της $f$ και το όριο $\lim\limits_{x \to -1^+} f(x)$. 
\item Να αποδείξετε ότι η $f$ είναι γνησίως αύξουσα στο πεδίο ορισμού της και να βρείτε το πρόσημό της σε αυτό.
\item Να αποδείξετε ότι $\ln(x+1) \ge \frac{x}{x+1}$ για κάθε $x \ge 0$.
\item Να βρείτε την εξίσωση της εφαπτομένης της $C_f$ στο σημείο με τετμημένη $x_0=0$.
\end{erwthma}
\end{thema}
%=========== ΤΕΛΟΣ - 22ο ΘΕΜΑ Β ===========

%=========== 23ο ΘΕΜΑ Β ===========
\begin{thema}{Β}
Δίνονται οι συναρτήσεις $g(x) = e^x - x$ και $f(x) = \ln(g(x))$.
\begin{erwthma}
\item Να μελετήσετε την $g$ ως προς τη μονοτονία, να βρείτε το ολικό της ελάχιστο και να αποδείξετε ότι το πεδίο ορισμού της $f$ είναι όλο το $\mathbb{R}$.
\item Να μελετήσετε την $f$ ως προς τη μονοτονία και τα ακρότατα.
\item Να αποδείξετε ότι $f(x) \ge 0$ για κάθε $x \in \mathbb{R}$.
\item Να υπολογίσετε το εμβαδόν του χωρίου που περικλείεται από τη γραφική παράσταση της $g$, τον άξονα $x'x$ και τις κατακόρυφες ευθείες $x=0$ και $x=1$.
\end{erwthma}
\end{thema}
%=========== ΤΕΛΟΣ - 23ο ΘΕΜΑ Β ===========

%=========== 24ο ΘΕΜΑ Β ===========
\begin{thema}{Β}
Δίνεται η συνάρτηση $f(x) = \sqrt{x^2+1} - x$.
\begin{erwthma}
\item Να βρείτε το πεδίο ορισμού της $f$ και να υπολογίσετε τα όρια $\lim\limits_{x \to -\infty} f(x)$ και $\lim\limits_{x \to +\infty} f(x)$.
\item Να αποδείξετε ότι η $f$ είναι γνησίως φθίνουσα και αυστηρά θετική στο $\mathbb{R}$.
\item Να αποδείξετε ότι η γραφική παράσταση της $f$ έχει ασύμπτωτες τις ευθείες $y = -2x$ στο $-\infty$ και $y = 0$ στο $+\infty$.
\item Να μελετήσετε την $f$ ως προς την κυρτότητα.
\end{erwthma}
\end{thema}
%=========== ΤΕΛΟΣ - 24ο ΘΕΜΑ Β ===========

%=========== 25ο ΘΕΜΑ Β ===========
\begin{thema}{Β}
Δίνεται η συνάρτηση $f(x) = \dfrac{\ln x}{x^2}$.
\begin{erwthma}
\item Να βρείτε το πεδίο ορισμού και τις ασύμπτωτες της γραφικής παράστασης της $f$.
\item Να μελετήσετε τη συνάρτηση $f$ ως προς τη μονοτονία και να βρείτε το ολικό ακρότατο.
\item Να αποδείξετε ότι $\ln{x} \le \dfrac{x^2}{2e}$ για κάθε $x > 0$.
\item Να υπολογίσετε το εμβαδόν του χωρίου που περικλείεται από τη γραφική παράσταση της $f$, τον άξονα $x'x$ και τις ευθείες $x=1$ και $x=e$.
\end{erwthma}
\end{thema}
%=========== ΤΕΛΟΣ - 25ο ΘΕΜΑ Β ===========

%=========== 26ο ΘΕΜΑ Β ===========
\begin{thema}{Β}
Δίνεται η συνάρτηση $f(x) = \dfrac{x-1}{x+1}$.
\begin{erwthma}
\item Να βρείτε το πεδίο ορισμού της $f$ και να αποδείξετε ότι αντιστρέφεται.
\item Να βρείτε το πεδίο ορισμού και τον τύπο της συνάρτησης $f^{-1}$.
\item Να βρείτε το πεδίο ορισμού της συνάρτησης $h = f \circ f$ και να αποδείξετε ότι $h(x) = -\frac{1}{x}$.
\item Να υπολογίσετε το όριο $\lim\limits_{x \to +\infty} \left( x \cdot f\left(\frac{1}{x}\right) \right)$.
\end{erwthma}
\end{thema}
%=========== ΤΕΛΟΣ - 26ο ΘΕΜΑ Β ===========

%=========== 27ο ΘΕΜΑ Β ===========
\begin{thema}{Β}
Δίνεται η συνάρτηση $f(x) = \dfrac{2x^2-1}{x^2+1}$.
\begin{erwthma}
\item Να βρείτε το πεδίο ορισμού και τις οριζόντιες ασύμπτωτες της γραφικής παράστασης της $f$.
\item Να μελετήσετε την $f$ ως προς τη μονοτονία και να βρείτε τα τοπικά ακρότατα.
\item Να βρείτε τα σημεία τομής της γραφικής παράστασης της $f$ με τους άξονες και την εξίσωση της εφαπτομένης της $C_f$ στο σημείο τομής της με τον άξονα $y'y$.
\item Να υπολογίσετε το εμβαδόν του χωρίου που περικλείεται από τη γραφική παράσταση της $f$, την ευθεία $y=2$ και τις ευθείες $x=0$, $x=1$.
\end{erwthma}
\end{thema}
%=========== ΤΕΛΟΣ - 27ο ΘΕΜΑ Β ===========

%=========== 28ο ΘΕΜΑ Β ===========
\begin{thema}{Β}
Δίνεται η συνάρτηση $f(x) = (x-1)e^{-x} + 2$.
\begin{erwthma}
\item Να μελετήσετε την $f$ ως προς τη μονοτονία και την κυρτότητα.
\item Να βρείτε το σημείο καμπής της $C_f$ καθώς και τα τοπικά ακρότατα της $f$.
\item Να αποδείξετε ότι η ευθεία $y = 2$ είναι οριζόντια ασύμπτωτη της $C_f$ στο $+\infty$.
\item Να υπολογίσετε το εμβαδόν του χωρίου που περικλείεται από τη γραφική παράσταση της $f$, τον άξονα $x'x$ και τις κατακόρυφες ευθείες $x=1$ και $x=2$.
\end{erwthma}
\end{thema}
%=========== ΤΕΛΟΣ - 28ο ΘΕΜΑ Β ===========

%=========== 29ο ΘΕΜΑ Β ===========
\begin{thema}{Β}
Δίνεται η συνάρτηση $f(x) = x \ln x - x + 1$.
\begin{erwthma}
\item Να βρείτε το πεδίο ορισμού της $f$ και το όριο $\lim\limits_{x \to 0^+} f(x)$.
\item Να μελετήσετε τη συνάρτηση $f$ ως προς τη μονοτονία και να βρείτε το ολικό της ελάχιστο.
\item Να αποδείξετε ότι ισχύει $x \ln x \ge x - 1$ για κάθε $x > 0$.
\item Να υπολογίσετε το εμβαδόν του χωρίου που περικλείεται από τη γραφική παράσταση της συνάρτησης $f$, τον άξονα $x'x$ και τις ευθείες $x=1$ και $x=e$.
\end{erwthma}
\end{thema}
%=========== ΤΕΛΟΣ - 29ο ΘΕΜΑ Β ===========

%=========== 30ο ΘΕΜΑ Β ===========
\begin{thema}{Β}
Δίνεται η συνάρτηση $f(x) = x^2 - 2\ln x$.
\begin{erwthma}
\item Να μελετήσετε την $f$ ως προς τη μονοτονία, τα ακρότατα και την κυρτότητα.
\item Να αποδείξετε ότι $x^2 \ge 2\ln x + 1$ για κάθε $x > 0$.
\item Να αποδείξετε ότι υπάρχει μοναδικό σημείο της $C_f$ στο οποίο η εφαπτομένη είναι παράλληλη στην ευθεία $\varepsilon: y = 2x - 3$ και να βρείτε την εξίσωση αυτής της εφαπτομένης.
\item Να αποδείξετε ότι η εξίσωση $f(x) = 2$ έχει τουλάχιστον μία πραγματική ρίζα στο διάστημα $(1, e)$.
\end{erwthma}
\end{thema}
%=========== ΤΕΛΟΣ - 30ο ΘΕΜΑ Β ===========

%=========== 31ο ΘΕΜΑ Β ===========
\begin{thema}{Β}
Δίνεται η συνάρτηση $f(x) = e^{2x} - 4e^x + 3$.
\begin{erwthma}
\item Να βρείτε τα σημεία τομής της γραφικής παράστασης της $f$ με τον άξονα $x'x$.
\item Να βρείτε τα όρια $\lim\limits_{x \to -\infty} f(x)$ και $\lim\limits_{x \to +\infty} f(x)$, και την οριζόντια ασύμπτωτη της $C_f$.
\item Να μελετήσετε την $f$ ως προς τη μονοτονία και να βρείτε το ολικό ελάχιστο.
\item Να υπολογίσετε το εμβαδόν του χωρίου που περικλείεται από τη $C_f$, τον άξονα $x'x$ και τις ευθείες $x=0$ και $x=\ln 3$.
\end{erwthma}
\end{thema}
%=========== ΤΕΛΟΣ - 31ο ΘΕΜΑ Β ===========

%=========== 32ο ΘΕΜΑ Β ===========
\begin{thema}{Β}
Δίνεται η συνάρτηση $f(x) = \dfrac{x^3}{x^2-1}$.
\begin{erwthma}
\item Να βρείτε το πεδίο ορισμού και τις κατακόρυφες ασύμπτωτες της γραφικής παράστασης της $f$.
\item Να αποδείξετε ότι η ευθεία $y = x$ είναι πλάγια ασύμπτωτη της $C_f$ στο $+\infty$ και στο $-\infty$.
\item Να μελετήσετε την $f$ ως προς τη μονοτονία και να βρείτε τις θέσεις τοπικών ακροτάτων.
\item Βρείτε το σύνολο τιμών της $f$ καθώς και το πλήθος ριζών της εξίσωσης $f(x)=\lambda$, για κάθε τιμή του πραγματικού αριθμού $\lambda\in\mathbb{R}$.
\end{erwthma}
\end{thema}
%=========== ΤΕΛΟΣ - 32ο ΘΕΜΑ Β ===========

%=========== 33ο ΘΕΜΑ Β ===========
\begin{thema}{Β}
Δίνεται η συνάρτηση $f(x) = \ln(e^x + 1) - x$.
\begin{erwthma}
\item Να αποδείξετε ότι $f(x) = \ln(1 + e^{-x})$ και να υπολογίσετε τα όρια $\lim\limits_{x \to -\infty} f(x)$ και $\lim\limits_{x \to +\infty} f(x)$.
\item Να βρείτε τις ασύμπτωτες της γραφικής παράστασης της $f$.
\item Να αποδείξετε ότι η $f$ αντιστρέφεται και να βρείτε τον τύπο της αντίστροφης συνάρτησης $f^{-1}$.
\item Να αποδείξετε ότι η εξίσωση $f(x) = x$ έχει μοναδική ρίζα στο διάστημα $(0, 1)$.
\end{erwthma}
\end{thema}
%=========== ΤΕΛΟΣ - 33ο ΘΕΜΑ Β ===========

%=========== 34ο ΘΕΜΑ Β ===========
\begin{thema}{Β}
Δίνεται η συνάρτηση $f(x) = x^2e^{-x+1}$.
\begin{erwthma}
\item Να μελετήσετε τη συνάρτηση $f$ ως προς τη μονοτονία και τα ακρότατα.
\item Να μελετήσετε τη συνάρτηση $f$ ως προς την κυρτότητα.
\item Να βρείτε την εξίσωση της εφαπτομένης της $C_f$ στο σημείο της $x_0 = 1$ και να δείξετε ότι $x^2e^{-x+1} \le x$ για κάθε $x \ge 0$.
\item Να υπολογίσετε το όριο $\lim\limits_{x \to 0^+} (f(x) \cdot \ln x)$.
\end{erwthma}
\end{thema}
%=========== ΤΕΛΟΣ - 34ο ΘΕΜΑ Β ===========

%=========== 35ο ΘΕΜΑ Β ===========
\begin{thema}{Β}
Δίνεται η συνάρτηση $f(x) = \dfrac{1+\ln x}{x}$.
\begin{erwthma}
\item Να βρείτε το πεδίο ορισμού και τις ασύμπτωτες της γραφικής παράστασης της $f$.
\item Να μελετήσετε τη συνάρτηση $f$ ως προς τη μονοτονία και να βρείτε το ολικό μέγιστό της.
\item Να αποδείξετε ότι $1+\ln x \le \dfrac{x}{e}$ για κάθε $x > 0$.
\item Να βρείτε το πλήθος των ριζών της εξίσωσης $1 + \ln x = \kappa x$ για τις διάφορες πραγματικές τιμές του $\kappa \in \mathbb{R}$.
\end{erwthma}
\end{thema}
%=========== ΤΕΛΟΣ - 35ο ΘΕΜΑ Β ===========

%=========== 36ο ΘΕΜΑ Β ===========
\begin{thema}{Β}
Δίνονται οι συναρτήσεις $f(x) = e^x - 1$ και $g(x) = \ln x$. 
\begin{erwthma}
\item Να βρείτε το πεδίο ορισμού και τον τύπο των συναρτήσεων $h = \dfrac{f}{g}$ και $k = f \cdot g$.
\item Να υπολογίσετε το όριο $\lim\limits_{x \to 1} h(x)$.
\item Να αποδείξετε ότι η εξίσωση $k(x) = 1$ έχει τουλάχιστον μία πραγματική ρίζα στο διάστημα $(1, e)$.
\item Να βρείτε τις ασύμπτωτες της γραφικής παράστασης της συνάρτησης $h$.
\end{erwthma}
\end{thema}
%=========== ΤΕΛΟΣ - 36ο ΘΕΜΑ Β ===========

%=========== 37ο ΘΕΜΑ Β ===========
\begin{thema}{Β}
Δίνεται η συνάρτηση $f(x) = x \ln x$ με πεδίο ορισμού $(0, +\infty)$.
\begin{erwthma}
\item Να υπολογίσετε το όριο $\lim\limits_{x \to 0^+} f(x)$ και στη συνέχεια, να βρείτε την κατακόρυφη ασύμπτωτη της $C_f$, αν υπάρχει.
\item Να μελετήσετε την $f$ ως προς τη μονοτονία και την κυρτότητα.
\item Να αποδείξετε ότι υπάρχει τουλάχιστον ένα $\xi \in (1, e)$ τέτοιο ώστε $f'(\xi) = \dfrac{e}{e-1}$.
\item Να υπολογίσετε το εμβαδόν του χωρίου που περικλείεται από την $C_f$, τον άξονα $x'x$ και τις ευθείες $x=1$ και $x=e$.
\end{erwthma}
\end{thema}
%=========== ΤΕΛΟΣ - 37ο ΘΕΜΑ Β ===========

%=========== 38ο ΘΕΜΑ Β ===========
\begin{thema}{Β}
Δίνεται η συνάρτηση $f(x) = x(x-2)e^x$ με πεδίο ορισμού το $\mathbb{R}$.
\begin{erwthma}
\item Να αποδείξετε ότι η εξίσωση $f(x) = 3$ έχει τουλάχιστον μία λύση στο $(2, 3)$.
\item Να υπολογίσετε το όριο $\lim\limits_{x \to -\infty} f(x)$.
\item Να μελετήσετε την $f$ ως προς τη μονοτονία και να βρείτε τα τοπικά ακρότατα.
\item Να βρείτε το σύνολο τιμών της $f$.
\end{erwthma}
\end{thema}
%=========== ΤΕΛΟΣ - 38ο ΘΕΜΑ Β ===========

%=========== 39ο ΘΕΜΑ Β ===========
\begin{thema}{Β}
Δίνεται η συνάρτηση $f(x) = \dfrac{e^x - x - 1}{x^2}$ για $x \ne 0$.
\begin{erwthma}
\item Να υπολογίσετε το όριο $\lim\limits_{x \to 0} f(x)$ και να βρείτε την οριζόντια ασύμπτωτη της $C_f$ στο $-\infty$.
\item Να ορίσετε τον αριθμό $a \in \mathbb{R}$ ώστε η συνάρτηση $g(x) = \begin{cases} f(x), & x \ne 0 \\ a, & x = 0 \end{cases}$ να είναι συνεχής στο $\mathbb{R}$.
\item Να αποδείξετε ότι υπάρχει $x_0 \in (0, 1)$ τέτοιο ώστε η συνεχής συνάρτηση $g$ να παίρνει την τιμή $g(x_0) = 0,7$.
\item Να αποδείξετε ότι $\lim\limits_{x \to +\infty} \dfrac{\ln x}{f(x)} = 0$.
\end{erwthma}
\end{thema}
%=========== ΤΕΛΟΣ - 39ο ΘΕΜΑ Β ===========

%=========== 40ο ΘΕΜΑ Β ===========
\begin{thema}{Β}
Δίνεται η συνάρτηση $f(x) = x^3 - 3x + 1$ με $x \in \mathbb{R}$.
\begin{erwthma}
\item Να μελετήσετε την $f$ ως προς τη μονοτονία.
\item Να αποδείξετε ότι η εξίσωση $f(x)=0$ έχει μοναδική ρίζα $x_0 \in (0, 1)$.
\item Να αποδείξετε ότι υπάρχει σημείο της γραφικής παράστασης της $f$ στο οποίο η εφαπτομένη είναι παράλληλη στο τμήμα που ενώνει τα σημεία $A(0, f(0))$ και $B(2, f(2))$ και να βρείτε την τετμημένη του $\xi \in (0, 2)$.
\item Βρείτε, άν υπάρχει, το όριο $\lim\limits_{x\to x_0}{\dfrac{\syn{x}}{f(x)}}$
\end{erwthma}
\end{thema}
%=========== ΤΕΛΟΣ - 40ο ΘΕΜΑ Β ===========

%=========== 41ο ΘΕΜΑ Β ===========
\begin{thema}{Β}
Δίνονται οι συναρτήσεις $f(x) = \ln(x^2+1)$ και $g(x) = x$, ορισμένες στο $\mathbb{R}$.
\begin{erwthma}
\item Να ορίσετε τη συνάρτηση $h(x) = f(x) - g(x)$ και να αποδείξετε ότι το όριο $\lim\limits_{x \to 0} \dfrac{h(x)}{x^2}$ ισούται με $1$.
\item Να μελετήσετε τη συνάρτηση $h$ ως προς τη μονοτονία.
\item Να μελετήσετε την $h$ ως προς την κυρτότητα και τα σημεία καμπής.
\item Να βρείτε, αν υπάρχει, την πλάγια ασύμπτωτη της γραφικής παράστασης της $h$ στο $+\infty$.
\end{erwthma}
\end{thema}
%=========== ΤΕΛΟΣ - 41ο ΘΕΜΑ Β ===========

%=========== 42ο ΘΕΜΑ Β ===========
\begin{thema}{Β}
Δίνεται η συνάρτηση $f(x) = \dfrac{\ln x}{x-1}$ για $x > 0$ και $x \ne 1$.
\begin{erwthma}
\item Να εξετάσετε αν η συνάρτηση
\[g(x)=\begin{cdcases}
f(x) &, 0<x\neq 1\\ 1 & , x=1
\end{cdcases}\]
είναι συνεχής.
\item Δίνεται η συνάρτηση $h(x) = (x-1)^2f(x)$. Να δείξετε ότι υπάρχει ρίζα της εξίσωσης $h(x) = 1$ στο διάστημα $(1, e)$.
\item Να βρείτε τις ασύμπτωτες της $C_f$.
\end{erwthma}
\end{thema}
%=========== ΤΕΛΟΣ - 42ο ΘΕΜΑ Β ===========

%=========== 43ο ΘΕΜΑ Β ===========
\begin{thema}{Β}
Δίνονται οι συναρτήσεις $f(x) = x+1$ και $g(x) = \dfrac{1}{x+1}$.
\begin{erwthma}
\item Να βρείτε το πεδίο ορισμού και τον τύπο της συνάρτησης $h = f + g$.
\item Να εξετάσετε αν η εξίσωση $h(x)=0$ έχει λύση στο διάστημα $[-2, 0]$ και αν υπάρχει $x_0 \in [0, 1]$ τέτοιο ώστε $h(x_0)=2$.
\item Να μελετήσετε την $h$ ως προς τη μονοτονία και να βρείτε τα τοπικά ακρότατα.
\item Να βρείτε τις ασύμπτωτες της γραφικής παράστασης της $h$.
\end{erwthma}
\end{thema}
%=========== ΤΕΛΟΣ - 43ο ΘΕΜΑ Β ===========

%=========== 44ο ΘΕΜΑ Β ===========
\begin{thema}{Β}
Δίνεται η συνάρτηση $f(x) = e^x - x - 1$ με $x \in \mathbb{R}$.
\begin{erwthma}
\item Να βρείτε το όριο $\lim\limits_{x \to 0} \dfrac{f(x)}{x^2}$.
\item Να αποδείξετε ότι $f(x) \ge 0$ για κάθε $x \in \mathbb{R}$.
\item Να εξετάσετε τη συνάρτηση ως προς την κυρτότητα.
\item Να δείξετε ότι η εξίσωση $f(x) = 2$ έχει ακριβώς μία ρίζα στο διάστημα $(1, 2)$.
\end{erwthma}
\end{thema}
%=========== ΤΕΛΟΣ - 44ο ΘΕΜΑ Β ===========

%=========== 45ο ΘΕΜΑ Β ===========
\begin{thema}{Β}
Δίνεται η συνάρτηση $f(x) = x e^x$ με πεδίο ορισμού το $\mathbb{R}$.
\begin{erwthma}
\item Να υπολογίσετε το $\lim\limits_{x \to -\infty} f(x)$ και να βρείτε ποια είναι η οριζόντια ασύμπτωτη στο $-\infty$.
\item Να βρείτε την εξίσωση της εφαπτομένης της $C_f$ που είναι παράλληλη με την ευθεία $\zeta: y=x+3$.
\item Να αποδείξετε ότι η γραφική παράσταση της $f$ τέμνει την ευθεία $y = -1$ σε ένα τουλάχιστον σημείο με την τετμημένη του στο διάστημα $(-1, 0)$.
\item Να προσδιορίσετε το ολικό ελάχιστο της συνάρτησης.
\end{erwthma}
\end{thema}
%=========== ΤΕΛΟΣ - 45ο ΘΕΜΑ Β ===========

%=========== 46ο ΘΕΜΑ Β ===========
\begin{thema}{Β}
Δίνεται η συνάρτηση $f(x) = \dfrac{\sqrt{x+1}-1}{x}$ για $x \ge -1, x \ne 0$.
\begin{erwthma}
\item Να υπολογίσετε το $\lim\limits_{x \to 0} f(x)$.
\item Έστω η συνεχής συνάρτηση
\[g(x)=\begin{cdcases}
f(x)&, x\in D_f\\ a-\frac{1}{2}&, x=0
\end{cdcases}\]
Να δείξετε ότι $a=1$.
\item Να αποδείξετε ότι η συνάρτηση $g$ λαμβάνει την τιμή $\dfrac{2}{5}$ για τουλάχιστον ένα $\xi \in (0, 3)$.
\item Να αποδείξετε ότι η $g$ αντιστρέφεται στο πεδίο ορισμού της.
\end{erwthma}
\end{thema}
%=========== ΤΕΛΟΣ - 46ο ΘΕΜΑ Β ===========

%=========== 47ο ΘΕΜΑ Β ===========
\begin{thema}{Β}
Δίνεται η πολυωνυμική συνάρτηση $f(x) = x^4 - 2x^2 + 1$.
\begin{erwthma}
\item Να αποδείξετε ότι υπάρχει τουλάχιστον ένα $\xi \in (-1, 1)$ στο οποίο η εφαπτομένη της γραφικής παράστασης της $f$ είναι οριζόντια.
\item Να μελετήσετε τη συνάρτηση ως προς τη μονοτονία και να εντοπίσετε τα τοπικά της ακρότατα.
\item Να βρείτε τα σημεία καμπής της $C_f$.
\item Υπολογίστε το εμβαδόν του χωρίου μεταξύ της $C_f$ και της ευθείας $y=9$.
\end{erwthma}
\end{thema}
%=========== ΤΕΛΟΣ - 47ο ΘΕΜΑ Β ===========

%=========== 48ο ΘΕΜΑ Β ===========
\begin{thema}{Β}
Δίνεται η συνάρτηση $f(x) = e^x (x-1) + 1$, με $x \in \mathbb{R}$.
\begin{erwthma}
\item Να υπολογίσετε το $\lim\limits_{x \to -\infty} f(x)$.
\item Να μελετήσετε την $f$ ως προς τη μονοτονία και την κυρτότητα.
\item Να αποδείξετε ότι υπάρχει μοναδικό $x_0 \in (0, 1)$ τέτοιο ώστε $f(x_0) = \dfrac{1}{2}$.
\item Να βρείτε την εξίσωση της εφαπτομένης της $C_f$ στο σημείο της $M(0, f(0))$.
\end{erwthma}
\end{thema}
%=========== ΤΕΛΟΣ - 48ο ΘΕΜΑ Β ===========

%=========== 49ο ΘΕΜΑ Β ===========
\begin{thema}{Β}
Δίνεται η συνάρτηση $f(x) = \dfrac{x - \hm x}{x + \hm x}$.
\begin{erwthma}
\item Να βρείτε το $\lim\limits_{x \to 0} f(x)$.
\item Να υπολογίσετε το $\lim\limits_{x \to +\infty} f(x)$.
\item Να αποδείξετε ότι η συνάρτηση $h(x) = x - \hm x$ είναι γνησίως αύξουσα.
\item Να δικαιολογήσετε γιατί η εξίσωση $x = \hm x$ έχει μοναδική ρίζα σε όλο το $\mathbb{R}$.
\end{erwthma}
\end{thema}
%=========== ΤΕΛΟΣ - 49ο ΘΕΜΑ Β ===========

%=========== 50ο ΘΕΜΑ Β ===========
\begin{thema}{Β}
Δίνεται η συνάρτηση $f(x) = \begin{cases} x^3 + ax + \beta, & x < 1 \\ \ln x + 2x, & x \ge 1 \end{cases}$.
\begin{erwthma}
  \item Να βρείτε τις τιμές των πραγματικών αριθμών $a, \beta$ ώστε η συνάρτηση να είναι παραγωγίσιμη στο $x_0=1$.
  \item Για $a = 0$ και $\beta = 1$, να προσδιορίσετε, αν υπάρχουν, τα κρίσιμα σημεία της συνάρτησης.
  \item Να αποδείξετε ότι η συνάρτηση $f$ είναι γνησίως αύξουσα σε όλο το $\mathbb{R}$ και να λύσετε την ανίσωση $f(x^2 - x) > 1$.
  \item Να υπολογίσετε το εμβαδόν του χωρίου που περικλείεται από τη γραφική παράσταση της $f$, τον άξονα $x'x$ και τις ευθείες $x=1, x=e$.
\end{erwthma}
\end{thema}
%=========== ΤΕΛΟΣ - 50ο ΘΕΜΑ Β ===========

%=========== 51ο ΘΕΜΑ Β ===========
\begin{thema}{Β}
Δίνεται η συνάρτηση $f(x) = e^{-x^2}$.
\begin{erwthma}
\item Να αποδείξετε ότι υπάρχει $\xi \in (-1, 1)$ στο οποίο η εφαπτομένη της γραφικής παράστασης της $f$ είναι οριζόντια και να το βρείτε.
\item Να μελετήσετε την $f$ ως προς τη μονοτονία και να βρείτε το μέγιστο της συνάρτησης.
\item Να βρείτε τα σημεία καμπής της γραφικής παράστασης της $f$.
\item Να υπολογίσετε το όριο $\lim\limits_{x \to +\infty} f(x)$ και να γράψετε την αντίστοιχη ασύμπτωτη.
\end{erwthma}
\end{thema}
%=========== ΤΕΛΟΣ - 51ο ΘΕΜΑ Β ===========

%=========== 52ο ΘΕΜΑ Β ===========
\begin{thema}{Β}
Δίνεται η συνάρτηση $f(x) = \dfrac{e^x}{x}$ για $x > 0$.
\begin{erwthma}
\item Να αποδείξετε ότι υπάρχει τουλάχιστον ένα $x_0 \in (1, 2)$ τέτοιο ώστε $3f(x_0) = e^2+e$.
\item Να ελέγξετε αν η εξίσωση $3f(x) = e^2+e$ έχει πάνω από μία λύση στο $(1, 2)$, μελετώντας τη μονοτονία της $f$.
\item Να βρείτε το $\lim\limits_{x \to +\infty} f(x)$.
\item Να βρείτε την εξίσωση της εφαπτομένης της $C_f$ η οποία διέρχεται από την αρχή των αξόνων $O(0,0)$.
\end{erwthma}
\end{thema}
%=========== ΤΕΛΟΣ - 52ο ΘΕΜΑ Β ===========

%=========== 53ο ΘΕΜΑ Β ===========
\begin{thema}{Β}
Δίνεται η συνάρτηση $f(x) = x^3 + x - 1$.
\begin{erwthma}
\item Να αποδείξετε ότι υπάρχει τουλάχιστον ένα $\xi \in (0, 1)$ τέτοιο ώστε η εφαπτομένη της $C_f$ στο $M(\xi, f(\xi))$ να έχει κλίση $1$. Βρείτε το ακριβές $\xi$.
\item Να αποδείξετε ότι η εξίσωση $f(x) = 0$ έχει τουλάχιστον μία ρίζα στο $(0, 1)$.
\item Να αποδείξετε ότι αυτή η ρίζα είναι μοναδική στο $\mathbb{R}$.
\item Υπολογίστε το εμβαδόν του χωρίου που περικλείεται από τη $C_f$, τον άξονα $x'x$ και τις ευθείες $x=1$, $x=2$.
\end{erwthma}
\end{thema}
%=========== ΤΕΛΟΣ - 53ο ΘΕΜΑ Β ===========

%=========== 54ο ΘΕΜΑ Β ===========
\begin{thema}{Β}
Δίνεται η συνάρτηση $f(x) = x^2 \ln x$ με $x > 0$.
\begin{erwthma}
\item Να υπολογίσετε το $\lim\limits_{x \to 0^+} f(x)$.
\item Εξετάστε αν η $C_f$ έχει οριζόντια εφαπτομένη στο σημείο $M(x_0,f(x_0))$ με $x_0\in(0,1)$.
\item Να αποδείξετε ότι υπάρχει $\xi \in (1, e)$ τέτοιο ώστε η εφαπτομένη της $C_f$ στο $\xi$ να είναι παράλληλη στο τμήμα με άκρα $A(1, f(1))$ και $B(e, f(e))$.
\item Να μελετήσετε την κυρτότητα της συνάρτησης.
\end{erwthma}
\end{thema}
%=========== ΤΕΛΟΣ - 54ο ΘΕΜΑ Β ===========

%=========== 55ο ΘΕΜΑ Β ===========
\begin{thema}{Β}
Δίνονται οι συναρτήσεις $f(x) = e^x + 2x$ και $g(x) = x^2 + 1$.
\begin{erwthma}
\item Να ορίσετε τον τύπο της συνάρτησης σύνθεσης $f \circ g$.
\item Να αποδείξετε ότι υπάρχει ένα τουλάχιστον $x_0 \in (0, 1)$ τέτοιο ώστε $f(x_0) = 3$.
\item Να αποδείξετε ότι η εξίσωση $f'(x) = 0$ δεν έχει λύση στο $\mathbb{R}$. Ποιο είναι το συμπέρασμα για τη μονοτονία της $f$;
\item Να δείξετε ότι υπάρχει $\xi \in (0, x_0)$ ώστε $f'(\xi) = \dfrac{2}{x_0}$.
\end{erwthma}
\end{thema}
%=========== ΤΕΛΟΣ - 55ο ΘΕΜΑ Β ===========

%=========== 56ο ΘΕΜΑ Β ===========
\begin{thema}{Β}
Δίνεται η συνάρτηση $f(x) = \dfrac{x^2+a}{x-\beta}$. Δίνεται ότι η γραφική παράσταση της $f$ διέρχεται από το σημείο $A(2, 5)$ και έχει κατακόρυφη ασύμπτωτη την ευθεία $x=1$.
\begin{erwthma}
  \item Να αποδείξετε ότι $a=1$ και $\beta=1$.
  \item Να βρείτε, αν υπάρχουν, τις οριζόντιες ή πλάγιες ασύμπτωτες της $C_f$.
  \item Να μελετήσετε την $f$ ως προς τη μονοτονία και να εντοπίσετε τα ακρότατα.
  \item Να υπολογίσετε το εμβαδόν του χωρίου που περικλείεται από την $C_f$, την πλάγια ασύμπτωτή της και τις ευθείες $x=2$ και $x=e+1$.
\end{erwthma}
\end{thema}
%=========== ΤΕΛΟΣ - 56ο ΘΕΜΑ Β ===========

%=========== 57ο ΘΕΜΑ Β ===========
\begin{thema}{Β}
Δίνεται η συνάρτηση $f(x) = a x^3 - \beta x^2 + 1$. Δίνεται ότι το όριο $\lim\limits_{x\to +\infty}\dfrac{f(x)}{x^3} = 2$ και ότι η $f$ παρουσιάζει τοπικό ακρότατο στο εσωτερικό σημείο $x_0 = 1$.
\begin{erwthma}
  \item Να αποδείξετε ότι $a=2$ και $\beta=3$.
  \item Να μελετήσετε τη συνάρτηση ως προς την κυρτότητα.
  \item Να αποδείξετε ότι η εξίσωση $f(x)=0$ έχει τουλάχιστον μία πραγματική ρίζα στο διάστημα $(-1, 0)$.
  \item Να βρείτε την εξίσωση της εφαπτομένης της $C_f$ η οποία είναι παράλληλη στην ευθεία $y = 12x-5$.
\end{erwthma}
\end{thema}
%=========== ΤΕΛΟΣ - 57ο ΘΕΜΑ Β ===========

%=========== 58ο ΘΕΜΑ Β ===========
\begin{thema}{Β}
Δίνεται η συνάρτηση $f(x) = \begin{cdcases} \dfrac{e^{ax}-1}{x}, & x \ne 0 \\ \beta, & x=0 \end{cdcases}$. Δίνεται ότι η γραφική παράσταση της $f$ διέρχεται από το σημείο $M(1, e-1)$ και ότι η $f$ είναι συνεχής στο $x_0 = 0$.
\begin{erwthma}
  \item Να αποδείξετε ότι $a=1$ και $\beta=1$.
  \item Να βρείτε την εξίσωση της εφαπτομένης της $C_f$ στο $x=1$.
  \item Να υπολογίσετε το $\lim\limits_{x \to -\infty} f(x)$.
  \item Να αποδείξετε ότι η συνάρτηση είναι παραγωγίσιμη στο $x_0 = 0$ και να βρείτε την τιμή της $f'(0)$.
\end{erwthma}
\end{thema}
%=========== ΤΕΛΟΣ - 58ο ΘΕΜΑ Β ===========

%=========== 59ο ΘΕΜΑ Β ===========
\begin{thema}{Β}
Δίνεται η συνάρτηση $f(x) = \begin{cdcases} a x^2+1, & x \le 1 \\ 2x+\beta, & x > 1 \end{cdcases}$. Δίνεται ότι η συνάρτηση $f$ είναι παραγωγίσιμη στο $x_0 = 1$.
\begin{erwthma}
  \item Να αποδείξετε ότι $a=1$ και $\beta=-1$.
  \item Να μελετήσετε την $f$ ως προς τη μονοτονία σε όλο το $\mathbb{R}$.
  \item Να αποδείξετε ότι υπάρχει μόνο ένα σημείο της $C_f$ στο οποίο η εφαπτομένη να είναι παράλληλη στην ευθεία $y = 4x+2$.
  \item Να υπολογίσετε το ολοκλήρωμα $\displaystyle\int_0^2 f(x)dx$.
\end{erwthma}
\end{thema}
%=========== ΤΕΛΟΣ - 59ο ΘΕΜΑ Β ===========

%=========== 60ο ΘΕΜΑ Β ===========
\begin{thema}{Β}
Δίνεται η συνάρτηση $f(x) = e^x + \beta$. Έστω ότι η ευθεία $y = \dfrac{x}{e}+\dfrac{2+e}{e}$ εφάπτεται της $C_f$ στο σημείο της με τετμημένη $x_0=-1$.
\begin{erwthma}
  \item Να αποδείξετε ότι $\beta=1$.
  \item Μελετήστε την $f$ ως προς την κυρτότητα και δείξτε ότι $e^{x+1}\geq x+2$ για κάθε $x\in\mathbb{R}$.
  \item Να βρείτε τις ασύμπτωτες της $C_f$ στο $-\infty$.
  \item Να αποδείξετε ότι η συνάρτηση $f$ αντιστρέφεται στο πεδίο ορισμού της και να υπολογίσετε το όριο $\lim\limits_{x \to 2} \dfrac{f^{-1}(x)}{x-2}$.
\end{erwthma}
\end{thema}
%=========== ΤΕΛΟΣ - 60ο ΘΕΜΑ Β ===========

%=========== 61ο ΘΕΜΑ Β ===========
\begin{thema}{Β}
Δίνονται οι συναρτήσεις $f(x) = \ln x$ και $g(x) = a x^2 + \beta$. Οι γραφικές τους παραστάσεις $C_f$ και $C_g$ δέχονται κοινή εφαπτομένη στο κοινό τους σημείο με τετμημένη $x_0=1$.
\begin{erwthma}
  \item Να αποδείξετε ότι $a=\dfrac{1}{2}$ και $\beta=-\dfrac{1}{2}$.
  \item Να ορίσετε τη συνάρτηση $h(x) = f(x) - g(x)$ και να τη μελετήσετε ως προς τη μονοτονία.
  \item Να αποδείξετε ότι $\ln x \le \dfrac{x^2-1}{2}$ για κάθε $x > 0$.
  \item Να υπολογίσετε το εμβαδόν του χωρίου που περικλείεται από τη γραφική παράσταση της $h(x)$, τον άξονα $x'x$ και τις ευθείες $x=1$ και $x=e$.
\end{erwthma}
\end{thema}
%=========== ΤΕΛΟΣ - 61ο ΘΕΜΑ Β ===========

%=========== 62ο ΘΕΜΑ Β ===========
\begin{thema}{Β}
Δίνεται η συνάρτηση $f(x) = x^3 + a x^2 + 3x + 1$. Θεωρούμε γνωστό ότι η $f$ είναι γνησίως αύξουσα στο $\mathbb{R}$ και ότι παρουσιάζει σημείο καμπής στο σημείο $A(1, f(1))$.
\begin{erwthma}
  \item Να αποδείξετε ότι $a=-3$. Κατόπιν, να επιβεβαιώσετε ότι για $a=-3$ η συνάρτηση είναι όντως γνησίως αύξουσα.
  \item Να δείξετε ότι υπάρχει ακριβώς ένα $x_0 \in (0, 1)$ τέτοιο ώστε $f(x_0)=2$.
  \item Να υπολογίσετε το όριο $\lim\limits_{x \to 0} \dfrac{f(x)-1}{x}$.
  \item Να μελετήσετε την $f$ ως προς την κυρτότητα και να γράψετε την εξίσωση της εφαπτομένης στο μοναδικό σημείο καμπής της.
\end{erwthma}
\end{thema}
%=========== ΤΕΛΟΣ - 62ο ΘΕΜΑ Β ===========

%=========== 63ο ΘΕΜΑ Β ===========
\begin{thema}{Β}
Δίνεται η συνάρτηση $f(x) = (x-a)e^x + \beta$. Η $f$ παρουσιάζει τοπικό ακρότατο στο εσωτερικό σημείο $x_0=1$ του πεδίου ορισμού της, στο οποίο η συνάρτηση παίρνει τιμή ίση με $2$.
\begin{erwthma}
  \item Να αποδείξετε ότι $a=2$ και $\beta = e+2$.
  \item Να μελετήσετε την $f$ ως προς τη μονοτονία και το είδος του ακροτάτου.
  \item Να αποδείξετε ότι ισχύει $e^x(x-2)+e \ge 0$ για κάθε $x \in \mathbb{R}$.
  \item Να βρείτε, εφόσον υπάρχουν, τις οριζόντιες ασύμπτωτες της γραφικής παράστασης της $f$.
\end{erwthma}
\end{thema}
%=========== ΤΕΛΟΣ - 63ο ΘΕΜΑ Β ===========

%=========== 64ο ΘΕΜΑ Β ===========
\begin{thema}{Β}
Δίνεται η συνάρτηση $f(x) = \ln(x^2+1) + a x^2 + \beta$. Η γραφική παράσταση της $f$ διέρχεται από την αρχή των αξόνων $O(0,0)$ και έχει σημείο καμπής το σημείο $M(1, \ln 2)$.
\begin{erwthma}
  \item Να αποδείξετε ότι $a=0$ και $\beta=0$.
  \item Να αποδείξετε ότι $\ln(x^2+1) \le x^2$ για κάθε $x \in \mathbb{R}$. Πότε ισχύει η ισότητα?
  \item Να υπολογίσετε το $\lim\limits_{x \to 0^+} \dfrac{f(x)}{x}$.
  \item Να αποδείξετε ότι η συνάρτηση παρουσιάζει ολικό ελάχιστο και να εξετάσετε αν η γραφική της παράσταση έχει σημεία καμπής.
\end{erwthma}
\end{thema}
%=========== ΤΕΛΟΣ - 64ο ΘΕΜΑ Β ===========

%=========== 65ο ΘΕΜΑ Β ===========
\begin{thema}{Β}
Δίνεται η συνάρτηση $f(x) = \dfrac{a e^x + \beta}{e^x + 1}$. Η $C_f$ έχει οριζόντιες ασύμπτωτες τις ευθείες $y=2$ στο $+\infty$ και $y=-1$ στο $-\infty$.
\begin{erwthma}
  \item Να αποδείξετε ότι $a=2$ και $\beta=-1$.
  \item Να αποδείξετε ότι η $f$ αντιστρέφεται στο πεδίο ορισμού της και να βρείτε την $f^{-1}$.
  \item Να υπολογίσετε το όριο $\lim\limits_{x \to 0} \dfrac{f(x)-\frac{1}{2}}{x}$.
  \item Να υπολογίσετε το ολοκλήρωμα $\displaystyle\int_0^{\ln 2} f(x)dx$.
\end{erwthma}
\end{thema}
%=========== ΤΕΛΟΣ - 65ο ΘΕΜΑ Β ===========

%=========== 66ο ΘΕΜΑ Β ===========
\begin{thema}{Β}
Δίνεται η συνάρτηση $f(x) = \dfrac{x^2+ax+\beta}{x-1}$. Έστω ότι η $C_f$ έχει πλάγια ασύμπτωτη στο $+\infty$ την ευθεία $y = x+2$ και η εφαπτομένη της στο σημείο με τετμημένη $0$ είναι οριζόντια.
\begin{erwthma}
  \item Να αποδείξετε ότι $a=1$ και $\beta=-1$.
  \item Να βρείτε την εξίσωση της κατακόρυφης ασύμπτωτης της συνάρτησης.
  \item 
  \item Να βρείτε, εφόσον υπάρχουν, τα τοπικά ακρότατα της συνάρτησης.
\end{erwthma}
\end{thema}
%=========== ΤΕΛΟΣ - 66ο ΘΕΜΑ Β ===========

%=========== 67ο ΘΕΜΑ Β ===========
\begin{thema}{Β}
Δίνεται η συνάρτηση $f(x) = a x^2 - x \ln x + \beta$, για $x>0$. Δίνεται ότι η $f$ παρουσιάζει ακρότατο στο σημείο $x_0=1$ με τιμή $\dfrac{3}{2}$.
\begin{erwthma}
  \item Να αποδείξετε ότι $a=\dfrac{1}{2}$ και $\beta=1$.
  \item Να μελετήσετε την κυρτότητα της $f$.
  \item Να αποδείξετε ότι το ακρότατο που εντοπίσατε είναι ολικό ελάχιστο.
  \item Να υπολογίσετε το όριο $\lim\limits_{x \to +\infty} \dfrac{f(x)}{x^2}$.
\end{erwthma}
\end{thema}
%=========== ΤΕΛΟΣ - 67ο ΘΕΜΑ Β ===========

%=========== 68ο ΘΕΜΑ Β ===========
\begin{thema}{Β}
Δίνεται η συνάρτηση $f(x) = \dfrac{x^2+a x}{x-2}$. Δίνεται ότι η ευθεία $x=2$ είναι κατακόρυφη ασύμπτωτη, και ότι η εφαπτομένη της $C_f$ στο σημείο με τετμημένη $x_0=1$ είναι παράλληλη στον άξονα $x'x$.
\begin{erwthma}
  \item Να αποδείξτε ότι $a=-4$.
  \item Να μελετήσετε την $f$ ως προς τη μονοτονία.
  \item Να αποδείξετε ότι υπάρχουν ακριβώς δύο σημεία της $C_f$ στα οποία η εφαπτομένη διέρχεται από το σημείο $M(2, 0)$.
  \item Να βρείτε, αν υπάρχουν, τις πλάγιες ασύμπτωτες της $C_f$ στο $+\infty$ και στο $-\infty$.
\end{erwthma}
\end{thema}
%=========== ΤΕΛΟΣ - 68ο ΘΕΜΑ Β ===========

%=========== 69ο ΘΕΜΑ Β ===========
\begin{thema}{Β}
Δίνεται η συνάρτηση $f(x) = \ln(x-a) + \beta x$. Δίνεται ότι η $f$ τέμνει τον άξονα $x'x$ στο σημείο με τετμημένη $2$ και η γραφική της παράσταση έχει κατακόρυφη ασύμπτωτη την ευθεία $x=1$.
\begin{erwthma}
  \item Να αποδείξετε ότι $a=1$ και $\beta=0$.
  \item Να βρείτε το πεδίο ορισμού και τη μονοτονία της $f$.
  \item Να βρείτε το πεδίο ορισμού και τον τύπο της αντίστροφης συνάρτησης $f^{-1}$.
  \item Να αποδείξετε ότι η εξίσωση $f(x) = \dfrac{1}{x-1}$ έχει ακριβώς μία ρίζα στο $(1, +\infty)$.
\end{erwthma}
\end{thema}
%=========== ΤΕΛΟΣ - 69ο ΘΕΜΑ Β ===========

%=========== 70ο ΘΕΜΑ Β ===========
\begin{thema}{Β}
Δίνεται η συνάρτηση $f(x) = e^{x-a} + \beta x$. Γνωρίζουμε ότι η $f$ παρουσιάζει ακρότατο στο σημείο $M(1, 0)$.
\begin{erwthma}
  \item Να αποδείξετε ότι $a=1$ και $\beta=-1$.
  \item Να αποδείξετε ότι το ακρότατο είναι ολικό ελάχιστο της συνάρτησης.
  \item Να αποδείξετε ότι $e^{x-1} \ge x$ για κάθε $x \in \mathbb{R}$.
  \item Να υπολογίσετε το εμβαδόν του χωρίου που περικλείεται από τη γραφική παράσταση της συνάρτησης, τον άξονα $x'x$ και τις ευθείες $x=0$ και $x=1$.
\end{erwthma}
\end{thema}
%=========== ΤΕΛΟΣ - 70ο ΘΕΜΑ Β ===========

%=========== 71ο ΘΕΜΑ Β ===========
\begin{thema}{Β}
Δίνεται η συνάρτηση $f(x) = a x e^x + \beta$. Δίνεται ότι το όριο $\lim\limits_{x \to -\infty} f(x) = 2$ και η $C_f$ διέρχεται από σημείο $A(1, e+2)$.
\begin{erwthma}
  \item Αποδείξτε ότι $\beta=2$. Στη συνέχεια, αποδείξτε ότι $a=1$.
  \item Να μελετήσετε τη συνάρτηση $f$ ως προς την κυρτότητα και να βρείτε το σημείο καμπής.
  \item Να αποδείξετε ότι η εξίσωση $f(x) = 3$ έχει τουλάχιστον μία πραγματική ρίζα στο διάστημα $(0, 1)$.
  \item Να υπολογίσετε το εμβαδόν της επιφάνειας που περικλείεται από την $C_f$, τους άξονες $x'x, y'y$ και την ευθεία $x=1$.
\end{erwthma}
\end{thema}
%=========== ΤΕΛΟΣ - 71ο ΘΕΜΑ Β ===========

%=========== 72ο ΘΕΜΑ Β ===========
\begin{thema}{Β}
Δίνεται η συνάρτηση $f(x) = x^3 + a x + \beta$. Η ευθεία $y = 3x - 4$ εφάπτεται στη γραφική παράσταση της $f$ στο σημείο της με τετμημένη $x_0=1$.
\begin{erwthma}
  \item Να εξηγήσετε γιατί $f(1) = -1$ και $f'(1) = 3$. Από αυτό, να αποδείξετε ότι $a=0$ και $\beta=-2$.
  \item Να αποδείξετε ότι υπάρχει σημείο $\xi \in (0, 2)$ στο οποίο η εφαπτομένη της $C_f$ να είναι παράλληλη στην ευθεία $y = 4x+1$.
  \item Να υπολογίσετε το όριο $\lim\limits_{x \to 2} \dfrac{f(x)-6}{x-2}$.
  \item Να αποδείξετε ότι η συνάρτηση είναι αντιστρέψιμη και να λύσετε την ανίσωση $f(f(x)) > f(6)$.
\end{erwthma}
\end{thema}
%=========== ΤΕΛΟΣ - 72ο ΘΕΜΑ Β ===========

%=========== 73ο ΘΕΜΑ Β ===========
\begin{thema}{Β}
Δίνεται η συνάρτηση $f(x) = \begin{cdcases} \dfrac{\ln x}{x-1}, & x\in(0,1) \\ a, & x=1 \end{cdcases}$. Γνωρίζουμε ότι η συνάρτηση είναι συνεχής στο σημείο $x_0=1$.
\begin{erwthma}
  \item Να αποδείξετε ότι $a=1$.
  \item Να βρείτε, αν υπάρχει, την κατακόρυφη ασύμπτωτη της συνάρτησης.
  \item Να εξετάσετε αν η $f$ είναι παραγωγίσιμη στο $x_0=1$.
  \item Να αποδείξετε ότι η συνάρτηση $f$ είναι γνησίως φθίνουσα στο διάστημα $(0, 1)$.
\end{erwthma}
\end{thema}
%=========== ΤΕΛΟΣ - 73ο ΘΕΜΑ Β ===========

%=========== 74ο ΘΕΜΑ Β ===========
\begin{thema}{Β}
Δίνεται η συνάρτηση $f(x) = \dfrac{x+a}{\beta x-2}$. Η $C_f$ έχει κατακόρυφη ασύμπτωτη την ευθεία $x=2$ και οριζόντια ασύμπτωτη την ευθεία $y=1$.
\begin{erwthma}
  \item Να αποδείξετε ότι $\beta=1$. Στη συνέχεια, βρείτε την τιμή του $a$, γνωρίζοντας επιπλέον ότι το σημείο $M(-2, 0) \in C_f$.
  \item Να αποδείξετε ότι η $f$ αντιστρέφεται στο πεδίο ορισμού της και να βρείτε τον τύπο της αντίστροφης συνάρτησης.
  \item Να αποδείξετε ότι η εξίσωση $f(x) = 2x$ έχει ακριβώς δύο λύσεις.
  \item Να υπολογίσετε το εμβαδόν του χωρίου μεταξύ της $C_f$, των αξόνων $x'x, y'y$, και της ευθείας $x=1$.
\end{erwthma}
\end{thema}
%=========== ΤΕΛΟΣ - 74ο ΘΕΜΑ Β ===========

%=========== 75ο ΘΕΜΑ Β ===========
\begin{thema}{Β}
Δίνεται η συνάρτηση $f(x) = \dfrac{ax^2+1}{x}$. Έστω ότι η $f$ παρουσιάζει τοπικό ακρότατο στο $x_0=1$.
\begin{erwthma}
  \item Να αποδείξετε ότι $a=1$.
  \item Να βρείτε τα τοπικά ακρότατα της $f$. Έχει ολικά ακρότατα η $f$?
  \item Να αποδείξετε ότι υπάρχει τουλάχιστον ένα $\xi \in (1, 2)$ τέτοιο ώστε $f'(\xi) = \dfrac{1}{2}$.
  \item Να βρείτε αν υπάρχουν τις ασύμπτωτες της $C_f$.
\end{erwthma}
\end{thema}
%=========== ΤΕΛΟΣ - 75ο ΘΕΜΑ Β ===========

%=========== 76ο ΘΕΜΑ Β ===========
\begin{thema}{Β}
Έστω ένα ορθογώνιο παράθυρο με ένα ημικύκλιο ενσωματωμένο στην πάνω πλευρά του.
\begin{center}
\begin{tikzpicture}[scale=0.8]
    \draw[thick, fill=maincolor!10] (0, 0) -- (4, 0) -- (4, 3) arc (0:180:2) -- cycle;
    \draw[thick, dashed] (0, 3) -- (4, 3);
    \draw[<->, >=latex] (0, -0.3) -- (4, -0.3) node[midway, below] {$2x$};
    \draw[<->, >=latex] (-0.3, 0) -- (-0.3, 3) node[midway, left] {$y$};
    \draw[<->, >=latex] (4.3, 0) -- (4.3, 3) node[midway, right] {$y$};
    \filldraw (2, 3) circle (1.5pt) node[below] {};
    \draw[->, >=latex] (2, 3) -- (2.8, 4.83) node[midway, right] {$x$};
\end{tikzpicture}
\end{center}
Το συνολικό μήκος του πλαισίου του παραθύρου (δηλαδή οι τρεις ευθύγραμμες πλευρές του ορθογωνίου και το τόξο του ημικυκλίου, χωρίς την ενδιάμεση βάση) είναι $10$ μέτρα. Θεωρούμε το πλάτος του παραθύρου ίσο με $2x$ και το ύψος του ορθογώνιου μέρους ίσο με $y$, όπου $x, y > 0$.
\begin{erwthma}
  \item Να αποδείξετε ότι το εμβαδόν ολόκληρου του παραθύρου (το εμβαδόν του ορθογωνίου συν το εμβαδόν του ημικυκλίου) ως συνάρτηση του $x$ δίνεται από τον τύπο $E(x) = 10x - 2x^2 - \frac{\pi}{2}x^2$ για $x \in \left(0, \frac{10}{2+\pi}\right)$.
  \item Να μελετήσετε την $E(x)$ ως προς τη μονοτονία και να βρείτε την τιμή του $x$ για την οποία το εμβαδόν μεγιστοποιείται. Τι παρατηρείτε για τη σχέση του $x$ με το $y$ στη θέση αυτή?
  \item Να βρείτε το σύνολο τιμών της $E(x)$.
  \item Με βάση το σύνολο τιμών, να εξετάσετε πόσα διαφορετικά παράθυρα μπορούν να κατασκευαστούν που να έχουν συνολικό εμβαδόν $5 \; \tau.\mu.$.
\end{erwthma}
\end{thema}
%=========== ΤΕΛΟΣ - 76ο ΘΕΜΑ Β ===========

%=========== 77ο ΘΕΜΑ Β ===========
\begin{thema}{Β}
Δίνεται το σταθερό σημείο $A(2, 0)$ και ένα σημείο $M$ που κινείται στη γραφική παράσταση της συνάρτησης $g(x) = \sqrt{x}$ με $x \ge 0$.
\begin{erwthma}
  \item Να αποδείξετε ότι η απόσταση $d(x) = (AM)$ δίνεται από τη σχέση $d(x) = \sqrt{x^2 - 3x + 4}$ για κάθε $x \ge 0$.
  \item Να μελετήσετε τη συνάρτηση $d(x)$ ως προς τη μονοτονία και να προσδιορίσετε τις συντεταγμένες του σημείου $M_{0}$ της καμπύλης που βρίσκεται πιο κοντά στο σημείο $A$.
  \item Να αποδείξετε ότι η συνάρτηση $d(x)$ έχει πλάγια ασύμπτωτη στο $+\infty$, την οποία και να βρείτε.
  \item Να υπολογίσετε το εμβαδόν του χωρίου που περικλείεται από την $g(x)$, τον άξονα $x'x$ και την κατακόρυφη ευθεία που διέρχεται από το $M_0$.
\end{erwthma}
\end{thema}
%=========== ΤΕΛΟΣ - 77ο ΘΕΜΑ Β ===========

%=========== 78ο ΘΕΜΑ Β ===========
\begin{thema}{Β}
Δίνεται μια συνάρτηση $f: \mathbb{R} \to \mathbb{R}$ για την οποία ισχύει $1 - x^2 \le f(x) \le 1 + x^2$ για κάθε $x \in \mathbb{R}$.
\begin{erwthma}
  \item Να υπολογίσετε το όριο $\lim\limits_{x \to 0} f(x)$. Αν επιπλέον δίνεται ότι η $f$ είναι συνεχής στο σημείο $x_0=0$, ποια είναι η τιμή $f(0)$;
  \item Να αποδείξετε ότι η συνάρτηση $f$ είναι παραγωγίσιμη στο σημείο $x_0=0$ και να βρείτε την τιμή της παραγώγου $f'(0)$.
  \item Να γράψετε, με βάση τα παραπάνω, την εξίσωση της εφαπτομένης της $C_f$ στο $x_0=0$.
  \item Να υπολογίσετε το όριο $\lim\limits_{x \to 0} \left[ f(x) \cdot \hm\left(\frac{1}{x}\right) \right]$.
\end{erwthma}
\end{thema}
%=========== ΤΕΛΟΣ - 78ο ΘΕΜΑ Β ===========

%=========== 79ο ΘΕΜΑ Β ===========
\begin{thema}{Β}
Δίνεται μια συνεχής συνάρτηση $f: (0, +\infty) \to \mathbb{R}$ για την οποία ικανοποιείται η ανισότητα \[\dfrac{x^2 - x}{x+1} \le f(x) \le x - 2 + \dfrac{2}{x}\] για κάθε $x > 0$.
\begin{erwthma}
  \item Να υπολογίσετε το όριο $\lim\limits_{x \to +\infty} \dfrac{f(x)}{x}$.
  \item Να υπολογίσετε το όριο $\lim\limits_{x \to +\infty} (f(x) - x)$.
  \item Να γράψετε ποια είναι η πλάγια ασύμπτωτη της $C_f$ στο $+\infty$ αιτιολογώντας την απάντησή σας.
  \item Αν επιπλέον ισχύει $\lim\limits_{x \to 1} \dfrac{f(x)}{x-1} = 2$, να βρείτε ποια ευθεία εφάπτεται στη γραφική παράσταση της $f$ στο σημείο με τετμημένη $x_0=1$.
\end{erwthma}
\end{thema}
%=========== ΤΕΛΟΣ - 79ο ΘΕΜΑ Β ===========

%=========== 80ο ΘΕΜΑ Β ===========
\begin{thema}{Β}
Δίνεται μια συνεχής συνάρτηση $f: \mathbb{R} \to \mathbb{R}$ για την οποία ικανοποιείται η συνθήκη $f(x-2) = (x-2)e^{3-x}$ για κάθε  $x\in\mathbb{R}$.
\begin{erwthma}
  \item Να αποδείξετε ότι ο τύπος της συνάρτησης είναι $f(x) = x e^{1-x}$ για κάθε $x \in \mathbb{R}$.
  \item Να μελετήσετε την $f$ ως προς τη μονοτονία, να προσδιορίσετε τη θέση και το είδος των ακροτάτων της.
  \item Να αποδείξετε ότι η γραφική παράσταση της συνάρτησης παρουσιάζει ακριβώς ένα σημείο καμπής, του οποίου να βρείτε τις συντεταγμένες.
  \item Να αποδείξετε ότι η εξίσωση $f(x) = 2$ είναι αδύνατη στο $\mathbb{R}$.
\end{erwthma}
\end{thema}
%=========== ΤΕΛΟΣ - 80ο ΘΕΜΑ Β ===========

%=========== 81ο ΘΕΜΑ Β ===========
\begin{thema}{Β}
Δίνεται συνεχής συνάρτηση $f: (0, +\infty) \to \mathbb{R}$ τέτοια ώστε να ισχύει $f(e^x) = x e^{-x} + 2$ για κάθε $x \in \mathbb{R}$.
\begin{erwthma}
  \item Να αποδείξετε ότι ο τύπος της συνάρτησης είναι $f(x) = \dfrac{\ln x}{x} + 2$ για $x > 0$.
  \item Να βρείτε τη μονοτονία και το συνολικό σύνολο τιμών της $f$.
  \item Να βρείτε το πλήθος των ριζών της εξίσωσης $\ln x =  (\lambda - 2) x$ για τις διάφορες πραγματικές τιμές της παραμέτρου $\lambda \in \mathbb{R}$.
  \item Να υπολογίσετε το όριο $\lim\limits_{x \to 0^+} [x^2 f(x)]$.
\end{erwthma}
\end{thema}
%=========== ΤΕΛΟΣ - 81ο ΘΕΜΑ Β ===========

%=========== 82ο ΘΕΜΑ Β ===========
\begin{thema}{Β}
Δίνεται η συνάρτηση $f(x) = x^3 - 3x^2 + 1$ με πεδίο ορισμού το $\mathbb{R}$.
\begin{erwthma}
  \item Να μελετήσετε την $f$ ως προς τη μονοτονία και τα ακρότατα.
  \item Να προσδιορίσετε το σύνολο τιμών της $f$. Έχει η συνάρτηση ολικό ακρότατο?
  \item Να βρείτε το πλήθος των ριζών της εξίσωσης $x^3 - 3x^2 + \lambda = 0$ για τις διάφορες τιμές του πραγματικού αριθμού $\lambda$.
  \item Να βρείτε τα σημεία της $C_f$ στα οποία η εφαπτομένη της είναι οριζόντια ευθεία (παράλληλη στον άξονα $x'x$).
\end{erwthma}
\end{thema}
%=========== ΤΕΛΟΣ - 82ο ΘΕΜΑ Β ===========

%=========== 83ο ΘΕΜΑ Β ===========
\begin{thema}{Β}
Δίνεται η συνάρτηση $f(x) = e^x - \lambda x$, όπου $\lambda > 0$.
\begin{erwthma}
  \item Να μελετήσετε τη μονοτονία της $f$ και να βρείτε σε ποιο σημείο $x_0$ παρουσιάζει ακρότατο.
  \item Να αποδείξετε ότι η ελάχιστη τιμή που περιλαμβάνει το σύνολο τιμών της συνάρτησης είναι ίση με $\lambda - \lambda \ln \lambda$.
  \item Για $\lambda = e$, να αποδείξετε ότι η εξίσωση $f(x) = 0$ έχει μοναδική ρίζα στο $\mathbb{R}$. Ποια είναι αυτή?
  \item Να υπολογίσετε το όριο $\lim\limits_{x \to +\infty} \dfrac{f(x)}{x^2}$.
\end{erwthma}
\end{thema}
%=========== ΤΕΛΟΣ - 83ο ΘΕΜΑ Β ===========

%=========== 84ο ΘΕΜΑ Β ===========
\begin{thema}{Β}
Δίνεται η παραγωγίσιμη συνάρτηση $f: \mathbb{R} \to \mathbb{R}$ η οποία ικανοποιεί τη σχέση $f(x) + e^{f(x)} = x + 1$ για κάθε $x \in \mathbb{R}$.
\begin{erwthma}
  \item Δείξτε ότι η συνάρτηση $g(x) = x + e^x$ είναι $1-1$ και ότι $f(0) = 0$.
  \item Να αποδείξετε ότι $f'(x) > 0$ για κάθε $x \in \mathbb{R}$.
  \item Να λύσετε την ανίσωση $f(x^2 - x) > 0$.
  \item Να αποδείξετε ότι υπάρχει ακριβώς ένα $x_0 \in (0, 1)$ τέτοιο ώστε $f(x_0) = \dfrac{1}{2}$.
\end{erwthma}
\end{thema}
%=========== ΤΕΛΟΣ - 84ο ΘΕΜΑ Β ===========

%=========== 85ο ΘΕΜΑ Β ===========
\begin{thema}{Β}
Έστω η συνάρτηση $f(x) = -x^2+4$ με $x \ge 0$. Μια τυχαία εφαπτομένη της γραφικής της παράστασης στο σημείο $M(x_0, f(x_0))$ με $x_0 \in (0, 2)$ τέμνει τους θετικούς ημιάξονες $Ox$ και $Oy$ στα σημεία $A$ και $B$ αντίστοιχα.
\begin{erwthma}
  \item Βρείτε την εξίσωση της εφαπτομένης και προσδιορίστε τις συντεταγμένες των σημείων τομής $A$ και $B$ συναρτήσει του $x_0$.
  \item Να αποδείξετε ότι το εμβαδόν του ορθογωνίου τριγώνου $OAB$ δίνεται από τον τύπο $E(x_0) = \dfrac{(x_0^2+4)^2}{4x_0}$.
  \item Να βρείτε το συγκεκριμένο $x_0 \in (0, 2)$ για το οποίο το σχηματιζόμενο τρίγωνο $OAB$ αποκτά το ελάχιστο δυνατό εμβαδόν.
  \item Να υπολογίσετε το ελάχιστο αυτό εμβαδόν.
\end{erwthma}
\end{thema}
%=========== ΤΕΛΟΣ - 85ο ΘΕΜΑ Β ===========

%=========== 86ο ΘΕΜΑ Β ===========
\begin{thema}{Β}
Δίνεται η συνάρτηση $f: \mathbb{R} \to \mathbb{R}$, η οποία για κάθε $x \ne 0$ ικανοποιεί την ανισότητα $\left| x f(x) - \hm x \right| \le x^2$. Υποθέτουμε επιπλέον ότι η $f$ είναι συνεχής στο σημείο $x_0=0$.
\begin{erwthma}
  \item Να αποδείξετε ότι $ -x+\dfrac{\hm{x}}{x}\le f(x)\leq x+\dfrac{\hm{x}}{x}$ για κάθε $x\neq 0$.
  \item Να βρείτε το $f(0)$.
  \item Να αποδείξετε με τον ορισμό ότι η συνάρτηση $f$ είναι παραγωγίσιμη στο $0$ και να βρείτε την $f'(0)$.
  \item Αν η $f(x)$ είναι άρτια συνάρτηση στο $\mathbb{R}$, τι σημαίνει αυτό για την τιμή $f'(0)$?
\end{erwthma}
\end{thema}
%=========== ΤΕΛΟΣ - 86ο ΘΕΜΑ Β ===========

%=========== 87ο ΘΕΜΑ Β ===========
\begin{thema}{Β}
Δίνεται μια συνάρτηση $f$ με πεδίο ορισμού το $\mathbb{R}$ η οποία ικανοποιεί τη συναρτησιακή σχέση $2f(x) + f(-x) = 3x^2 + x$ για κάθε $x \in \mathbb{R}$.
\begin{erwthma}
  \item Δείξτε ότι ο τύπος της είναι $f(x) = x^2 + x$.
  \item Να μελετήσετε την $f(x)$ ως προς τη μονοτονία και την κυρτότητα.
  \item Να αποδείξετε ότι το σημείο $Α(-1, 0)$ ανήκει στη $C_f$ και ότι η εφαπτομένη σε αυτό διέρχεται από το σημείο $Β(0, -1)$.
  \item Να υπολογίσετε το εμβαδόν του χωρίου που περικλείεται από την $C_f$, τον άξονα $x'x$ και την παραπάνω εφαπτομένη.
\end{erwthma}
\end{thema}
%=========== ΤΕΛΟΣ - 87ο ΘΕΜΑ Β ===========

%=========== 88ο ΘΕΜΑ Β ===========
\begin{thema}{Β}
Δίνεται η συνάρτηση $f(x) = x^2 - 4x + 3$. Η γραφική της παράσταση τέμνει τον $x'x$ στα σημεία με τετμημένες $1$ και $3$ και παρουσιάζει ολικό ελάχιστο. Θεωρούμε επίσης τη συνάρτηση $g(x) = \ln(f(x))$.
\begin{erwthma}
  \item Βρείτε το πεδίο ορισμού της $g(x)$.
  \item Να υπολογίσετε τα πλευρικά όρια $\lim\limits_{x \to 1^-} g(x)$ και $\lim\limits_{x \to 3^+} g(x)$. Ποια είναι η γεωμετρική τους ερμηνεία?
  \item Να μελετήσετε τη συνάρτηση $g$ ως προς τη μονοτονία.
  \item Να αποδείξετε ότι η εξίσωση $g(x) = 0$ έχει ακριβώς δύο λύσεις.
\end{erwthma}
\end{thema}
%=========== ΤΕΛΟΣ - 88ο ΘΕΜΑ Β ===========

%===============================================================
\subsection*{Ασκήσεις Οπτικής Επίλυσης}
%=========== 89ο ΘΕΜΑ Β ===========
\begin{thema}{Β}
Στο παρακάτω σχήμα δίνεται η γραφική παράσταση μιας συνεχούς συνάρτησης $f$.
\begin{center}
\begin{tikzpicture}
    \begin{axis}[axis lines = middle, axis line style = {->, >=latex}, grid = major, grid style = {very thin, gray!30}, xlabel = {$x$}, ylabel = {$y$}, xtick = {-2,1,2,3,4}, ytick = {1,2,3,4}, every axis x label/.style = {at={(current axis.right of origin)}, anchor=west}, every axis y label/.style = {at={(current axis.above origin)}, anchor=south}, width = 7cm, height = 6cm, samples = 100, xmin=-3.5, xmax=5.5, ymin=-0.5, ymax=4.5, xtick distance=1, ytick distance=1]
        \addplot[very thick, maincolor, domain=-3.5:1.75] {1/((x-2)^2)};
        \addplot[very thick, maincolor, domain=2.25:5.5] {1 + 1/((x-2)^2)};
        \addplot[dashed, thick] coordinates {(2,-0.5) (2,4.5)};
        \addplot[dashed, thick] coordinates {(2.25,1) (5.5,1)};
        \node[below left] at (axis cs:0,0) {$0$};
    \end{axis}
\end{tikzpicture}
\end{center}
Με βάση το σχήμα, απαντήστε στα εξής:
\begin{erwthma}
  \item Να βρείτε το πεδίο ορισμού της $f$ καθώς και τα όρια $\lim\limits_{x \to 2^-}f(x)$, $\lim\limits_{x \to 2^+}f(x)$, $\lim\limits_{x \to -\infty}f(x)$ και $\lim\limits_{x \to +\infty}f(x)$.
  \item Να υπολογίσετε το όριο $\lim\limits_{x \to 2} e^{\frac{1}{f(x)}}$.
  \item Να υπολογίσετε το όριο $\lim\limits_{x \to 1} \ln(f(x))$.
  \item Προσδιορίστε, αν υπάρχουν, τοπικά μέγιστα ή ελάχιστα σημεία της $f$ δικαιολογώντας την απάντησή σας.
\end{erwthma}
\end{thema}
%=========== ΤΕΛΟΣ - 89ο ΘΕΜΑ Β ===========

%=========== 90ο ΘΕΜΑ Β ===========
\begin{thema}{Β}
Δίνεται η γραφική παράσταση μιας συνάρτησης $g(x)$.
\begin{center}
\begin{tikzpicture}
    \begin{axis}[axis lines = middle, axis line style = {->, >=latex}, grid = major, grid style = {very thin, gray!30}, xlabel = {$x$}, ylabel = {$y$}, xtick = {-4,-3,-2,-1,1,2,3,4,5,6,7}, ytick = {-1,1,2}, every axis x label/.style = {at={(current axis.right of origin)}, anchor=west}, every axis y label/.style = {at={(current axis.above origin)}, anchor=south}, width = 9cm, height = 7cm, xmin=-4.5, xmax=7, ymin=-1.5, ymax=2.5, xtick distance=1, ytick distance=1]
        \addplot[very thick, \xrwma, smooth, tension=0.5] coordinates {(-4.5,0.05) (-3.5, 0.1) (-2.5, 0.5) (-1, 2) (0, 0) (1, -1) (2, 0) (3, 1.2) (5.5, 1.85) (7,1.98)};
        \addplot[dashed, thick] coordinates {(-1,0) (-1,2) (0,2)};
        \addplot[dashed, thick] coordinates {(1,0) (1,-1) (0,-1)};
        \addplot[dashed, thick] coordinates {(0,2) (6.5,2)};
    \end{axis}
\end{tikzpicture}
\end{center}

\begin{erwthma}
  \item Να βρείτε από το σχήμα τα όρια $\lim\limits_{x \to +\infty} g(x)$ και $\lim\limits_{x \to -\infty} g(x)$.
  \item Να υπολογίσετε το όριο $\lim\limits_{x \to -\infty} g(g(x))$.
  \item Να βρείτε, εφόσον υπάρχει, το όριο $\lim\limits_{x \to 0} \dfrac{1}{g(x)}$.
\end{erwthma}
\end{thema}
%=========== ΤΕΛΟΣ - 90ο ΘΕΜΑ Β ===========

%=========== 91ο ΘΕΜΑ Β ===========
\begin{thema}{Β}
Στο παρακάτω σχήμα δίνεται η γραφική παράσταση μιας συνάρτησης $f$.
\begin{center}
\begin{tikzpicture}
    \begin{axis}[axis lines = middle,
        axis line style = {->, >=latex},
        grid = major,
        grid style = {very thin, gray!30},
        xlabel = {$x$},
        ylabel = {$y$},
        xtick = {-3,-2,-1,1,2,3,4},
        ytick = {-2,-1,1,2,3,4,5},
        every axis x label/.style = {at={(current axis.right of origin)}, anchor=west},
        every axis y label/.style = {at={(current axis.above origin)}, anchor=south},
        width = 8cm, height = 7cm, samples = 100,
        xmin=-3.5, xmax=4.5, ymin=-2.5, ymax=6,
        xtick distance=1, ytick distance=1
    ]
        \addplot[very thick, maincolor, domain=1.15:4.5] {2 + 1/(x-1)};
        \addplot[very thick, maincolor, domain=-3.5:0.85] {2 + 1/(x-1)};
        \addplot[dashed, thick] coordinates {(1,-2.5) (1,6)};
        \addplot[dashed, thick] coordinates {(-3.5,2) (4.5,2)};
        \node[below left] at (axis cs:0,0) {$0$};
    \end{axis}
\end{tikzpicture}
\end{center}
Με βάση τη γραφική παράσταση της $C_f$:
\begin{erwthma}
  \item Να βρείτε το πεδίο ορισμού της συνάρτησης $f$ και να υπολογίσετε τα όρια $\lim\limits_{x \to 1^-} f(x)$, $\lim\limits_{x \to 1^+} f(x)$, $\lim\limits_{x \to -\infty} f(x)$ και $\lim\limits_{x \to +\infty} f(x)$.
  \item Να γράψετε τις εξισώσεις των ασύμπτωτων της γραφικής παράστασης της $f$.
  \item Αν γνωρίζετε ότι ο τύπος της συνάρτησης είναι της μορφής $f(x) = \dfrac{ax+\beta}{x-\gamma}$, να αποδείξετε, με τη βοήθεια της γραφικής παράστασης, ότι $f(x) = \dfrac{2x-1}{x-1}$.
  \item Να αποδείξετε ότι η συνάρτηση $f$ αντιστρέφεται, να βρείτε την αντίστροφη συνάρτηση $f^{-1}$ και να υπολογίσετε το όριο $\lim\limits_{x \to 1} \dfrac{f^{-1}(x) + x^2 - 1}{x-1}$.
\end{erwthma}
\end{thema}
%=========== ΤΕΛΟΣ - 91ο ΘΕΜΑ Β ===========

%=========== 92ο ΘΕΜΑ Β ===========
\begin{thema}{Β}
Στο παρακάτω σχήμα βλέπουμε τη γραφική παράσταση της \textbf{{παραγώγου}} $f'$ μιας συνάρτησης $f$, η οποία είναι παραγωγίσιμη στο $\mathbb{R}$.
\begin{center}
\begin{tikzpicture}
    \begin{axis}[axis lines = middle,
        axis line style = {->, >=latex},
        grid = major,
        grid style = {very thin, gray!30},
        xlabel = {$x$},
        ylabel = {$y$},
        xtick = {-2,-1,1,2,3,4,5},
        ytick = {-4,-3,-2,-1,1,2,3},
        every axis x label/.style = {at={(current axis.right of origin)}, anchor=west},
        every axis y label/.style = {at={(current axis.above origin)}, anchor=south},
        width = 8cm, height = 7cm, samples = 100,
        xmin=-2.5, xmax=4.5, ymin=-4.5, ymax=3,
        xtick distance=1, ytick distance=1
    ]
        \addplot[very thick, \xrwma, domain=-1.5:3.5] {x^2 - 2*x - 3};
        \node[below left] at (axis cs:0,0) {$0$};
    \end{axis}
\end{tikzpicture}
\end{center}
Με βάση το σχήμα:
\begin{erwthma}
  \item Να μελετήσετε τη συνάρτηση $f$ ως προς τη μονοτονία.
  \item Να βρείτε τις θέσεις και το είδος των τοπικών ακροτάτων της $f$.
  \item Να μελετήσετε τη συνάρτηση $f$ ως προς την κυρτότητα και τα σημεία καμπής.
  \item Αν επιπλέον γνωρίζετε ότι $f(0) = 2$, να βρείτε την εξίσωση της εφαπτομένης της $C_f$ στο σημείο της με τετμημένη $x_0=0$.
\end{erwthma}
\end{thema}
%=========== ΤΕΛΟΣ - 92ο ΘΕΜΑ Β ===========

%=========== 93ο ΘΕΜΑ Β ===========
\begin{thema}{Β}
Στο σχήμα δίνεται η γραφική παράσταση της \textbf{{δεύτερης παραγώγου}} $f''$ μιας συνάρτησης $f$, παραγωγίσιμης δύο φορές στο $\mathbb{R}$.
\begin{center}
\begin{tikzpicture}
    \begin{axis}[axis lines = middle,
        axis line style = {->, >=latex},
        grid = major,
        grid style = {very thin, gray!30},
        xlabel = {$x$},
        ylabel = {$y$},
        xtick = {-2,-1,1,2,3,4},
        ytick = {-3,-2,-1,1,2,3},
        every axis x label/.style = {at={(current axis.right of origin)}, anchor=west},
        every axis y label/.style = {at={(current axis.above origin)}, anchor=south},
        width = 8cm, height = 6cm, samples = 100,
        xmin=-2.5, xmax=4.5, ymin=-3.5, ymax=3.5,
        xtick distance=1, ytick distance=1
    ]
        \addplot[very thick, \xrwma, domain=-1.5:3.5] {x - 1};
        \node[below left] at (axis cs:0,0) {$0$};
    \end{axis}
\end{tikzpicture}
\end{center}
\begin{erwthma}
  \item Να μελετήσετε την $f$ ως προς την κυρτότητα και να βρείτε το σημείο καμπής της $C_f$.
  \item Αν είναι γνωστό ότι η $f$ παρουσιάζει τοπικό ακρότατο στη θέση $x_1=0$, να βρείτε αν πρόκειται για τοπικό μέγιστο ή τοπικό ελάχιστο, αιτιολογώντας την απάντησή σας.
  \item Με βάση το προηγούμενο ερώτημα και αν γνωρίζετε ότι $f'(2)=0$, να υπολογίσετε το ολοκλήρωμα $\displaystyle\int_0^2 f''(x)dx$ χωρίς να βρείτε τον τύπο της $f''(x)$.
\end{erwthma}
\end{thema}
%=========== ΤΕΛΟΣ - 93ο ΘΕΜΑ Β ===========

%=========== 94ο ΘΕΜΑ Β ===========
\begin{thema}{Β}
Στο παρακάτω σχήμα δίνεται η γραφική παράσταση της \textbf{{παραγώγου}} $f'$ μιας συνεχούς συνάρτησης $f:[0, 4] \to \mathbb{R}$.
\begin{center}
\begin{tikzpicture}
    \begin{axis}[axis lines = middle,
        axis line style = {->, >=latex},
        grid = major,
        grid style = {very thin, gray!30},
        xlabel = {$x$},
        ylabel = {$y$},
        xtick = {1,2,3,4},
        ytick = {-1,1,2,3},
        every axis x label/.style = {at={(current axis.right of origin)}, anchor=west},
        every axis y label/.style = {at={(current axis.above origin)}, anchor=south},
        width = 7cm, height = 5cm, samples = 100,
        xmin=-0.5, xmax=4.5, ymin=-1.5, ymax=3.5,
        xtick distance=1, ytick distance=1
    ]
        \addplot[very thick, maincolor, domain=0:2] {x};
        \addplot[very thick, maincolor, domain=2:4] {4-x};
        \node[below left] at (axis cs:0,0) {$0$};
    \end{axis}
\end{tikzpicture}
\end{center}
\begin{erwthma}
  \item Να μελετήσετε την $f$ ως προς τη μονοτονία στο $[0, 4]$.
  \item Να βρείτε τα ακρότατα της $f$.
  \item Η συνάρτηση $f$ έχει σημεία καμπής? Δικαιολογήστε την απάντησή σας.
  \item Να υπολογίσετε το ολοκλήρωμα $\displaystyle\int_0^4 f'(x)dx$ γεωμετρικά. Τι εκφράζει αυτό το αποτέλεσμα για τις τιμές της $f$?
\end{erwthma}
\end{thema}
%=========== ΤΕΛΟΣ - 94ο ΘΕΜΑ Β ===========

%=========== 95ο ΘΕΜΑ Β ===========
\begin{thema}{Β}
Δίνεται η γραφική παράσταση της συνάρτησης $f$ και η εφαπτομένη $(\epsilon)$ στο σημείο $M(0, f(0))$.
\begin{center}
\begin{tikzpicture}
    \begin{axis}[axis lines = middle,
        axis line style = {->, >=latex},
        grid = major,
        grid style = {very thin, gray!30},
        xlabel = {$x$},
        ylabel = {$y$},
        xtick = {-2,-1,1,2,3},
        ytick = {-1,1,2,3,4},
        every axis x label/.style = {at={(current axis.right of origin)}, anchor=west},
        every axis y label/.style = {at={(current axis.above origin)}, anchor=south},
        width = 8cm, height = 6cm, samples = 100,
        xmin=-2.5, xmax=3.5, ymin=-1.5, ymax=5,
        xtick distance=1, ytick distance=1
    ]
        \addplot[very thick, \xrwma, domain=-2:2] {exp(x)};
        \addplot[black, thick, domain=-1:3] {x+1};
        \filldraw[black] (axis cs:0,1) circle (1.5pt) node[anchor=north west] {$M(0,1)$};
        \filldraw[black] (axis cs:-1,0) circle (1.5pt) node[anchor=south east] {$A(-1,0)$};
        \node[below left] at (axis cs:0,0) {$0$};
    \end{axis}
\end{tikzpicture}
\end{center}
\begin{erwthma}
  \item Να βρείτε την εξίσωση της εφαπτομένης $(\epsilon)$.
  \item Να βρείτε τις τιμές $f(0)$ και $f'(0)$.
  \item Να υπολογίσετε το όριο $\lim\limits_{x \to 0} \dfrac{f(x)-1}{x}$. Στη συνέχεια επιβεβαιώστε το εφαρμόζοντας τον κανόνα \eng{De L'Hospital}.
  \item Παρατηρώντας τη γεωμετρική σχέση μεταξύ της $C_f$ και της $(\epsilon)$, να συμπεράνετε το πρόσημο της δεύτερης παραγώγου της $f$.
\end{erwthma}
\end{thema}
%=========== ΤΕΛΟΣ - 95ο ΘΕΜΑ Β ===========

%=========== 96ο ΘΕΜΑ Β ===========
\begin{thema}{Β}
Δίνονται οι συναρτήσεις $f(x) = \sqrt{x}$ και $g(x) = x^2$ με πεδίο ορισμού το $[0, +\infty)$. Στο ίδιο σύστημα συντεταγμένων δίνονται οι γραφικές τους παραστάσεις $C_f$ και $C_g$.
\begin{center}
\begin{tikzpicture}
    \begin{axis}[axis lines = middle,
        axis line style = {->, >=latex},
        grid = major,
        grid style = {very thin, gray!30},
        xlabel = {$x$},
        ylabel = {$y$},
        xtick = {1,2,3},
        ytick = {1,2,3},
        every axis x label/.style = {at={(current axis.right of origin)}, anchor=west},
        every axis y label/.style = {at={(current axis.above origin)}, anchor=south},
        width = 6cm, height = 6cm, samples = 100,
        xmin=-0.5, xmax=2.5, ymin=-0.5, ymax=2.5,
        xtick distance=1, ytick distance=1
    ]
        \addplot[very thick, \xrwma, domain=0:2.2] {sqrt(x)} node[pos=0.8,below] {$C_f$};
        \addplot[very thick, maincolor, domain=-0.5:1.6] {x^2} node[pos=0.8,right]{$C_g$};
        \node[below left] at (axis cs:0,0) {$0$};
    \end{axis}
\end{tikzpicture}
\end{center}
\begin{erwthma}
  \item Να αποδείξετε αλγεβρικά ότι η συνάρτηση $g$ αντιστρέφεται στο $[0, +\infty)$ με $g^{-1}(x) = f(x)$ και να λύσετε την εξίσωση $f(x) = g(x)$. Πώς επιβεβαιώνονται τα παραπάνω συμπεράσματα παρατηρώντας τις $C_f$ και $C_g$;
  \item Να βρείτε την εξίσωση της εφαπτομένης της $C_g$ στο σημείο της με τετμημένη $x_0=1$. Να αποδείξετε ότι $g(x) \ge 2x - 1$ για κάθε $x \ge 0$.
  \item Έστω συνάρτηση $h: (0, 1] \to \mathbb{R}$ για την οποία ικανοποιείται η σχέση $g(x) \hm\left(\dfrac{1}{x}\right) \le h(x) \le f(x)$ για κάθε $x \in (0, 1]$. Να υπολογίσετε το όριο $\lim\limits_{x \to 0^+} h(x)$.
  \item Να υπολογίσετε το εμβαδόν του χωρίου που περικλείεται από τις δύο καμπύλες.
\end{erwthma}
\end{thema}
%=========== ΤΕΛΟΣ - 96ο ΘΕΜΑ Β ===========

%=========== 97ο ΘΕΜΑ Β ===========
\begin{thema}{Β}
Δίνεται η γραφική παράσταση της συνάρτησης $f:[0, 3] \to \mathbb{R}$.
\begin{center}
\begin{tikzpicture}
    \begin{axis}[axis lines = middle,
        axis line style = {->, >=latex},
        grid = major,
        grid style = {very thin, gray!30},
        xlabel = {$x$},
        ylabel = {$y$},
        xtick = {1,2,3},
        ytick = {1,2,3},
        every axis x label/.style = {at={(current axis.right of origin)}, anchor=west},
        every axis y label/.style = {at={(current axis.above origin)}, anchor=south},
        width = 6cm, height = 6cm, samples = 100,
        xmin=-0.5, xmax=3.5, ymin=-0.5, ymax=3.5,
        xtick distance=1, ytick distance=1
    ]
        \addplot[very thick, \xrwma, domain=0:1] {x^2};
        \addplot[very thick, \xrwma, domain=1:3] {x};
        \node[below left] at (axis cs:0,0) {$0$};
        \filldraw[\xrwma] (axis cs:3,3) circle (1.5pt);
    \end{axis}
\end{tikzpicture}
\end{center}
\begin{erwthma}
  \item Από τη γραφική παράσταση, να βρείτε το σύνολο τιμών της συνάρτησης, τα ολικά της ακρότατα και να αιτιολογήσετε γεωμετρικά αν είναι παραγωγίσιμη ή όχι στο σημείο $x_0=1$.
  \item Αν δίνεται πλέον ότι ο τύπος της συνάρτησης είναι $f(x) = \begin{cases} x^2, & 0 \le x \le 1 \\ x, & 1 < x \le 3 \end{cases}$, να επιβεβαιώσετε αναλυτικά την απάντηση που δώσατε στο προηγούμενο ερώτημα για την παραγωγισιμότητα στο $1$.
  \item Υπολογίστε το ολοκλήρωμα $\displaystyle\int_0^2 f(x)dx$.
\end{erwthma}
\end{thema}
%=========== ΤΕΛΟΣ - 97ο ΘΕΜΑ Β ===========

%=========== 98ο ΘΕΜΑ Β ===========
\begin{thema}{Β}
Δίνεται η συνάρτηση $f(x) = \left| x-2 \right| - 1$ και η γραφική της παράσταση, όπως φαίνεται στο παρακάτω σχήμα.
\begin{center}
\begin{tikzpicture}
    \begin{axis}[axis lines = middle,
        axis line style = {->, >=latex},
        grid = major,
        grid style = {very thin, gray!30},
        xlabel = {$x$},
        ylabel = {$y$},
        xtick = {-1,1,2,3,4},
        ytick = {-1,1,2},
        every axis x label/.style = {at={(current axis.right of origin)}, anchor=west},
        every axis y label/.style = {at={(current axis.above origin)}, anchor=south},
        width = 8cm, height = 5cm, samples = 100,
        xmin=-0.5, xmax=4.5, ymin=-1.5, ymax=2.5,
        xtick distance=1, ytick distance=1
    ]
        \addplot[very thick, maincolor, domain=0:4] {abs(x-2)-1};
        \node[below left] at (axis cs:0,0) {$0$};
    \end{axis}
\end{tikzpicture}
\end{center}
\begin{erwthma}
  \item Να μελετήσετε, με τη βοήθεια του σχήματος, τη συνάρτηση $f$ ως προς τη μονοτονία και να βρείτε το σύνολο τιμών της. Στη συνέχεια, να βρείτε το πλήθος των ριζών της εξίσωσης $f(x) = \lambda$ για τις διάφορες τιμές της παραμέτρου $\lambda \in \mathbb{R}$.
  \item Να εξετάσετε γραφικά αν η $C_f$ δέχεται εφαπτομένη στο σημείο με τετμημένη $x_0=2$. Ποια είναι η εξίσωση της εφαπτομένης της $C_f$ στο σημείο με τετμημένη $x_1=4$;
  \item Να υπολογίσετε τα όρια $\lim\limits_{x \to 1} \dfrac{f(x)}{x-1}$ και $\lim\limits_{x \to 3} \dfrac{f(x)}{x-3}$. Ποια γεωμετρική έννοια της γραφικής παράστασης εκφράζουν οι τιμές αυτών των ορίων;
  \item Να γράψετε το ολοκλήρωμα που εκφράζει το εμβαδόν του χωρίου που περικλείεται από την $C_f$ και τον άξονα $x'x$. Πώς μπορείτε να υπολογίσετε το αποτέλεσμα αυτού του ολοκληρώματος χρησιμοποιώντας αποκλειστικά στοιχειώδη γεωμετρία;
\end{erwthma}
\end{thema}
%=========== ΤΕΛΟΣ - 98ο ΘΕΜΑ Β ===========

%=========== 99ο ΘΕΜΑ Β ===========
\begin{thema}{Β}
Δίνεται η γραφική παράσταση μιας συνάρτησης $f$.
\begin{center}
\begin{tikzpicture}
    \begin{axis}[axis lines = middle,
        axis line style = {->, >=latex},
        grid = major,
        grid style = {very thin, gray!30},
        xlabel = {$x$},
        ylabel = {$y$},
        xtick = {-1,1,2,3},
        ytick = {-3,-2,-1,1,2},
        every axis x label/.style = {at={(current axis.right of origin)}, anchor=west},
        every axis y label/.style = {at={(current axis.above origin)}, anchor=south},
        width = 7cm, height = 6cm, samples = 100,
        xmin=-1.5, xmax=3.5, ymin=-3.5, ymax=2.5,
        xtick distance=1, ytick distance=1
    ]
        \addplot[very thick, maincolor, domain=-2:3.5] {x^3 - 3*x^2 + 1};
        \node[below left] at (axis cs:0,0) {$0$};
    \end{axis}
\end{tikzpicture}
\end{center}
\begin{erwthma}
  \item Βρείτε το πρόσημο της συνάρτησης $f'$ και να εξηγήστε γιατί η εξίσωση $f'(x)=0$ έχει τουλάχιστον δύο ρίζες.
  \item Γνωρίζοντας ότι η $f$ είναι πολυωνυμική εξετάστε αν η συνάρτηση $f'$ ικανοποιεί τις υποθέσεις του θεωρήματος \eng{\eng{Rolle}} στο διάστημα $[0,2]$
  \item Δικαιολογήστε γιατί υπάρχει τουλάχιστον ένα $\xi \in (0, 2)$ τέτοιο ώστε η εφαπτομένη της $C_f$ να είναι παράλληλη στην ευθεία που διέρχεται από τα σημεία $A(1, 1)$ και $B(-2, 7)$.
\end{erwthma}
\end{thema}
%=========== ΤΕΛΟΣ - 99ο ΘΕΜΑ Β ===========

%=========== 100ο ΘΕΜΑ Β ===========
\begin{thema}{Β}
Στο παρακάτω σχήμα δίνεται η γραφική παράσταση της \textbf{{παραγώγου}} $f'$ μιας συνεχούς συνάρτησης $f(x)$ ορισμένης στο $(0, +\infty)$.
\begin{center}
\begin{tikzpicture}
    \begin{axis}[axis lines = middle,
        axis line style = {->, >=latex},
        grid = major,
        grid style = {very thin, gray!30},
        xlabel = {$x$},
        ylabel = {$y$},
        xtick = {1,2,3,4},
        ytick = {1,2,3,4,5},
        every axis x label/.style = {at={(current axis.right of origin)}, anchor=west},
        every axis y label/.style = {at={(current axis.above origin)}, anchor=south},
        width = 8cm, height = 6cm, samples = 100,
        xmin=-0.5, xmax=4.5, ymin=-0.5, ymax=5.5,
        xtick distance=1, ytick distance=1
    ]
        \addplot[very thick, \xrwma, domain=0.2:4] {1/x};
        \node[below left] at (axis cs:0,0) {$0$};
    \end{axis}
\end{tikzpicture}
\end{center}
\begin{erwthma}
  \item Παρατηρώντας το σχήμα, βγάλτε συμπεράσματα για τη μονοτονία και την κυρτότητα της $f(x)$.
  \item Να βρείτε, αν γνωρίζουμε ότι $f'(x) = \dfrac{1}{x}$ και $f(1)=0$, τον τύπο της συνάρτησης $f$.
  \item Να αποδείξετε ότι ισχύει $f(x) \le x-1$ για κάθε $x > 0$.
  \item Υπολογίστε το όριο $\lim\limits_{x \to 0^+} (x f(x))$.
\end{erwthma}
\end{thema}
%=========== ΤΕΛΟΣ - 100ο ΘΕΜΑ Β ===========

%=========== 101ο ΘΕΜΑ Β ===========
\begin{thema}{Β}
Στο παρακάτω σχήμα δίνεται η γραφική παράσταση $C_f$ μιας συνάρτησης $f$ ορισμένης στο $[-3, 4]$. Η $C_f$ αποτελείται από ευθύγραμμα τμήματα με κορυφές τα σημεία $A(-3, -1)$, $B(-1, 3)$, $\Gamma(1, -1)$ και $\Delta(4, 2)$.

\begin{center}
\begin{tikzpicture}[scale=0.85]
  % grid
  \draw[help lines, gray!30] (-4,-2) grid (5,4);
  % axes
  \draw[-latex] (-4.3,0) -- (5.3,0) node[right] {$x$};
  \draw[-latex] (0,-2.3) -- (0,4.3) node[above] {$y$};
  % ticks
  \foreach \x in {-3,-2,-1,1,2,3,4}
    \draw (\x,0.1) -- (\x,-0.1) node[below,font=\footnotesize] {$\x$};
  \foreach \y in {-1,1,2,3}
    \draw (0.1,\y) -- (-0.1,\y) node[left,font=\footnotesize] {$\y$};
  % curve
  \draw[very thick, maincolor] (-3,-1) -- (-1,3) -- (1,-1) -- (4,2);
  % points
  \filldraw[maincolor] (-3,-1) circle (2pt) node[below left] {$A$};
  \filldraw[maincolor] (-1,3) circle (2pt) node[above] {$B$};
  \filldraw[maincolor] (1,-1) circle (2pt) node[below] {$\Gamma$};
  \filldraw[maincolor] (4,2) circle (2pt) node[above right] {$\Delta$};
\end{tikzpicture}
\end{center}

\begin{erwthma}
  \item Βασιζόμενοι στα δεδομένα του σχήματος, να προσδιορίσετε τον τύπο της συνάρτησης $f$.
  \item Να μελετήσετε την εξίσωση $f(x) = \lambda$, με $\lambda \in \mathbb{R}$, ως προς το πλήθος των ριζών της. Στη συνέχεια, αν για $\lambda \in (-1, 2]$ η εξίσωση έχει ακριβώς τρεις διαφορετικές ρίζες $\rho_1, \rho_2, \rho_3$, να αποδείξετε ότι $\rho_1 + \rho_2 + \rho_3 = \lambda$.
  \item Σε ποια σημεία η $f$ δεν είναι παραγωγίσιμη; Να αιτιολογήσετε.
  \item Να υπολογίσετε γεωμετρικά το ολοκλήρωμα $\dint_{-3}^{4} |f(x)| \, dx$.
\end{erwthma}
\end{thema}
%=========== ΤΕΛΟΣ - 101ο ΘΕΜΑ Β ===========

%---------------------------------------------------------------

%=========== 102ο ΘΕΜΑ Β ===========
\begin{thema}{Β}
Στο παρακάτω σχήμα δίνεται η γραφική παράσταση $C_f$ μιας συνάρτησης $f$ ορισμένης στο $\mathbb{R} - \{1\}$.
\begin{center}
\begin{tikzpicture}[scale=0.8]
  % grid
  \draw[help lines, gray!30] (-3,-3) grid (7,5);
  % axes
  \draw[-latex] (-3.3,0) -- (7.3,0) node[right] {$x$};
  \draw[-latex] (0,-3.3) -- (0,5.3) node[above] {$y$};
  % ticks
  \foreach \x in {-2,-1,2,3,4,5,6}
    \draw (\x,0.1) -- (\x,-0.1) node[below,font=\footnotesize] {$\x$};
  \foreach \y in {-2,-1,1,2,3,4}
    \draw (0.1,\y) -- (-0.1,\y) node[left,font=\footnotesize] {$\y$};
  % asymptotes
  \draw[dashed, secondcolor] (1,-3.3) -- (1,5.3);
  \draw[dashed, secondcolor] (-3.3,2) -- (7.3,2);
  % left branch
  \draw[very thick, maincolor, domain=-3:0.75, samples=80] plot (\x, {2 + 1/(\x-1)});
  % right branch
  \draw[very thick, maincolor, domain=1.25:7, samples=80] plot (\x, {2 + 1/(\x-1)});
  % points
  \filldraw[maincolor] (0,1) circle (2pt) node[below left] {$(0,1)$};
  \filldraw[maincolor] (2,3) circle (2pt) node[above left] {$(2,3)$};
\end{tikzpicture}
\end{center}

\begin{erwthma}
  \item Παρατηρώντας το σχήμα, να γράψετε τις εξισώσεις των ασύμπτωτων ευθειών της $C_f$ και να γράψετε τα αντίστοιχα όρια που τις επιβεβαιώνουν.
  \item Αν γνωρίζετε ότι $f(x) = \dfrac{2x-1}{x-1}$ για κάθε $x \ne 1$, να αποδείξετε ότι καμία εφαπτομένη της $C_f$ δεν διέρχεται από το σημείο τομής των ασύμπτωτών της.
  \item Να εξετάσετε αν υπάρχει το όριο $\lim\limits_{x\to 1}{\dfrac{1}{f(x)-2}}$
  \item Να βρείτε το πεδίο ορισμού της συνάρτησης $f\circ f$.
\end{erwthma}
\end{thema}
%=========== ΤΕΛΟΣ - 102ο ΘΕΜΑ Β ===========

%---------------------------------------------------------------

%=========== 103ο ΘΕΜΑ Β ===========
\begin{thema}{Β}
Στο παρακάτω σχήμα δίνεται η γραφική παράσταση της \textbf{παραγώγου} $f'$ μιας συνεχούς συνάρτησης $f$ ορισμένης στο $[-2, 5]$.

\begin{center}
\begin{tikzpicture}[scale=0.75]
  % grid
  \draw[help lines, gray!30] (-3,-3) grid (6,5);
  % axes
  \draw[-latex] (-3.3,0) -- (6.3,0) node[right] {$x$};
  \draw[-latex] (0,-3.3) -- (0,5.3) node[above] {$y$};
  % ticks
  \foreach \x in {-2,-1,1,2,3,4,5}
    \draw (\x,0.1) -- (\x,-0.1) node[below,font=\footnotesize] {$\x$};
  \foreach \y in {-2,-1,1,2,3,4}
    \draw (0.1,\y) -- (-0.1,\y) node[left,font=\footnotesize] {$\y$};
  % f' curve (two line segments)
  \draw[very thick, secondcolor] (-2,-2) -- (1,4) -- (5,0);
  % shade positive area
  \fill[maincolor!10, opacity=0.5] (-2,-2) -- (-1,0) -- (-2,0) -- cycle;
  \fill[orange!10, opacity=0.5] (-1,0) -- (1,4) -- (5,0) -- (-1,0) -- cycle;
  % key points
  \filldraw[secondcolor] (-2,-2) circle (2pt);
  \filldraw[secondcolor] (1,4) circle (2pt) node[above right] {$(1,4)$};
  \filldraw[secondcolor] (5,0) circle (2pt);
  \filldraw[secondcolor] (-1,0) circle (2pt) node[above left,font=\footnotesize] {$(-1,0)$};
  % label
  \node[secondcolor] at (3.5,3) {$C_{f'}$};
\end{tikzpicture}
\end{center}

\begin{erwthma}
  \item Να μελετήσετε τη μονοτονία της $f$ και να βρείτε τα ακρότατά της.
  \item Να υπολογίσετε τη διαφορά $f(5) - f(-1)$ χρησιμοποιώντας το εμβαδόν του χρωματισμένου χωρίου.
  \item Αν $f(1) = 3$, να βρείτε την εξίσωση της εφαπτομένης της $C_f$ στο σημείο με τετμημένη $x_0 = 1$.
  \item Να μελετήσετε την κυρτότητα της $f$ στο $[-2, 5]$.
\end{erwthma}
\end{thema}
%=========== ΤΕΛΟΣ - 103ο ΘΕΜΑ Β ===========

%---------------------------------------------------------------
% Ασκήσεις 4-6: Περισσότερες Cf
%---------------------------------------------------------------

%=========== 104ο ΘΕΜΑ Β ===========
\begin{thema}{Β}
Στο παρακάτω σχήμα δίνεται η γραφική παράσταση $C_f$ μιας συνεχούς συνάρτησης $f$ ορισμένης στο $[-2, 3]$. Η $C_f$ αποτελείται από ένα τεταρτοκύκλιο με κέντρο την αρχή των αξόνων και ακτίνα $\rho=2$, στο διάστημα $[-2, 0]$ και ένα ευθύγραμμο τμήμα με άκρα τα σημεία $A(0, 2)$ και $B(3, -1)$.

\begin{center}
\begin{tikzpicture}[scale=0.85]
  \draw[help lines, gray!30] (-3,-2) grid (4,3);
  \draw[-latex] (-3.3,0) -- (4.3,0) node[right] {$x$};
  \draw[-latex] (0,-2.3) -- (0,3.3) node[above] {$y$};
  \foreach \x in {-2,-1,1,2,3}
    \draw (\x,0.1) -- (\x,-0.1) node[below,font=\footnotesize] {$\x$};
  \foreach \y in {-1,1,2}
    \draw (0.1,\y) -- (-0.1,\y) node[left,font=\footnotesize] {$\y$};
  % quarter circle
  \draw[very thick, maincolor!70!black] (-2,0) arc (180:90:2);
  % line segment
  \draw[very thick, maincolor!70!black] (0,2) -- (3,-1);
  \filldraw[maincolor!70!black] (-2,0) circle (2pt);
  \filldraw[maincolor!70!black] (0,2) circle (2pt) node[above right,font=\footnotesize] {$A(0,2)$};
  \filldraw[maincolor!70!black] (3,-1) circle (2pt) node[below right,font=\footnotesize] {$B(3,-1)$};
\end{tikzpicture}
\end{center}

\begin{erwthma}
  \item Βασιζόμενοι στα δεδομένα, να αποδείξετε ότι ο τύπος της συνάρτησης είναι $f(x) = \begin{cases} \sqrt{4-x^2}, & -2 \le x \le 0 \\ -x+2, & 0 < x \le 3 \end{cases}$.
  \item Να εξετάσετε αν η συνάρτηση $f$ είναι παραγωγίσιμη στο $x_0 = 0$. Τι συμπεραίνετε για την ύπαρξη εφαπτομένης της $C_f$ στο σημείο $A(0, 2)$;
  \item Να αιτιολογήσετε γραφικά γιατί η συνάρτηση $f$ δεν αντιστρέφεται στο $[-2, 3]$. Αν $g$ είναι ο περιορισμός της $f$ στο διάστημα $[-2, 0]$, να ορίσετε την αντίστροφη συνάρτηση $g^{-1}$.
  \item Να υπολογίσετε γεωμετρικά το ολοκλήρωμα $\dint_{0}^{3} f(x) \, dx$.
\end{erwthma}
\end{thema}
%=========== ΤΕΛΟΣ - 104ο ΘΕΜΑ Β ===========

%---------------------------------------------------------------

%=========== 105ο ΘΕΜΑ Β ===========
\begin{thema}{Β}
Στο παρακάτω σχήμα δίνεται η γραφική παράσταση μιας συνάρτησης $f$ ορισμένης στο $\mathbb{R}$, η οποία παρουσιάζει πλάγια ασύμπτωτη $y = \frac{1}{2}x$ στο $+\infty$ και οριζόντια ασύμπτωτη $y = -1$ στο $-\infty$.

\begin{center}
\begin{tikzpicture}[scale=0.65]
  \draw[help lines, gray!30] (-6,-3) grid (8,5);
  \draw[-latex] (-6.3,0) -- (8.3,0) node[right] {$x$};
  \draw[-latex] (0,-3.3) -- (0,5.3) node[above] {$y$};
  \foreach \x in {-5,-4,-3,-2,-1,1,2,3,4,5,6,7}
    \draw (\x,0.1) -- (\x,-0.1) node[below,font=\tiny] {$\x$};
  \foreach \y in {-2,-1,1,2,3,4}
    \draw (0.1,\y) -- (-0.1,\y) node[left,font=\footnotesize] {$\y$};
  
  % asymptotes
  \draw[dashed, thick] (-6,-1) -- (2,-1) node[below,font=\footnotesize] {$y=-1$};
  \draw[dashed, thick] (-1,-0.5) -- (8,4) node[right,font=\footnotesize] {$y=\frac{x}{2}$};
  
  % curve - strictly ONE draw command
  \draw[very thick, maincolor, smooth, tension=0.4]
    plot coordinates {(-6,-0.98) (-4,-0.9) (-2,-0.5) (0,0) (2,0.5) (4,1.6) (6,2.8) (8,3.9)};
    
  \filldraw[maincolor] (0,0) circle (2pt) node[above left,font=\footnotesize] {$O$};
\end{tikzpicture}
\end{center}

\begin{erwthma}
  \item Να βρείτε τα $\displaystyle\lim\limits_{x \to -\infty} f(x)$ και $\displaystyle\lim\limits_{x \to +\infty} \frac{f(x)}{x}$.
  \item Να υπολογίσετε το όριο $\displaystyle\lim\limits_{x \to +\infty} \frac{x f(x) - x^2 + 3}{2x^2 + x + 1}$.
  \item Αν γνωρίζετε ότι η συνάρτηση $f$ είναι γνησίως αύξουσα στο $\mathbb{R}$, να λύσετε την ανίσωση $f\left(|x-1| - 2\right) < 0$.
  \item Αν η $C_f$ βρίσκεται πάνω από την οριζόντια ασύμπτωτή της για κάθε $x \in \mathbb{R}$, να αποδείξετε ότι το εμβαδόν $E$ του χωρίου που περικλείεται από την $C_f$, τον άξονα $x'x$ και τις ευθείες $x=-2$ και $x=0$, είναι αυστηρά μικρότερο από $2$.
\end{erwthma}
\end{thema}
%=========== ΤΕΛΟΣ - 105ο ΘΕΜΑ Β ===========

%---------------------------------------------------------------

%=========== 106ο ΘΕΜΑ Β ===========
\begin{thema}{Β}
Στο παρακάτω σχήμα δίνεται η γραφική παράσταση $C_f$ μιας παραγωγίσιμης συνάρτησης $f$ με πεδίο ορισμού το διάστημα $[-2, 2]$.

\begin{center}
\begin{tikzpicture}[scale=1]
  \draw[help lines, gray!30] (-2.5,-2.5) grid (2.5,2.5);
  \draw[-latex] (-2.8,0) -- (2.8,0) node[right] {$x$};
  \draw[-latex] (0,-2.8) -- (0,2.8) node[above] {$y$};
  \foreach \x in {-2,-1,1,2}
    \draw (\x,0.1) -- (\x,-0.1) node[below,font=\footnotesize] {$\x$};
  \foreach \y in {-2,-1,1,2}
    \draw (0.1,\y) -- (-0.1,\y) node[left,font=\footnotesize] {$\y$};
  % curve
  \draw[very thick, maincolor!70!black, smooth, samples=100, domain=-2:2] plot (\x, {\x^3 - 3*\x});
  % points
  \filldraw[maincolor!70!black] (-1,2) circle (2pt) node[above] {$A(-1,2)$};
  \filldraw[maincolor!70!black] (1,-2) circle (2pt) node[below] {$B(1,-2)$};
  \filldraw[maincolor!70!black] (0,0) circle (2pt) node[below left] {$O$};
  \node[below right,font=\footnotesize] at (1.732,0) {$\rho_2$};
  \node[above left,font=\footnotesize] at (-1.732,0) {$\rho_1$};
\end{tikzpicture}
\end{center}

\begin{erwthma}
  \item Με βάση το σχήμα, να προσδιορίσετε τα διαστήματα στα οποία η συνάρτηση $f$ είναι κυρτή ή κοίλη, καθώς και τις συντεταγμένες του σημείου καμπής της.
  \item Να αποδείξετε ότι υπάρχει τουλάχιστον ένα $\xi \in (0, 2)$ τέτοιο ώστε η εφαπτομένη της $C_f$ στο σημείο $M(\xi, f(\xi))$ να είναι παράλληλη στην ευθεία $y = x$.
  \item Χρησιμοποιώντας αποκλειστικά το σχήμα, να αποδείξετε ότι $\dint_{0}^{1} f(x) \, dx > -2$.
  \item Να βρείτε το πλήθος των ριζών της εξίσωσης $f(f(x)) = 0$ στο διάστημα $[-2, 2]$.
\end{erwthma}
\end{thema}
%=========== ΤΕΛΟΣ - 106ο ΘΕΜΑ Β ===========

%---------------------------------------------------------------
% Ασκήσεις 7-9: Cf' reading
%---------------------------------------------------------------

%=========== 107ο ΘΕΜΑ Β ===========
\begin{thema}{Β}
Στο σχήμα δίνεται η $C_{f'}$ μιας παραγωγίσιμης $f:[-3, 4]\to\mathbb{R}$.
\begin{center}
\begin{tikzpicture}[scale=0.7]
  \draw[help lines, gray!30] (-4,-3) grid (5,4);
  \draw[-latex] (-4.3,0) -- (5.3,0) node[right] {$x$};
  \draw[-latex] (0,-3.3) -- (0,4.3) node[above] {$y$};
  \foreach \x in {-3,-2,-1,1,2,3,4}
    \draw (\x,0.1) -- (\x,-0.1) node[below,font=\footnotesize] {$\x$};
  \foreach \y in {-2,-1,1,2,3}
    \draw (0.1,\y) -- (-0.1,\y) node[left,font=\footnotesize] {$\y$};
  \draw[very thick, secondcolor, smooth, tension=0.5]
    plot coordinates {(-3,2) (-2,1.5) (-1,0) (0,-1.5) (0.5,-2) (1,-1.5) (2,0) (3,2) (4,3)};
  \filldraw[secondcolor] (-1,0) circle (2pt) node[above,font=\footnotesize] {$(-1,0)$};
  \filldraw[secondcolor] (2,0) circle (2pt) node[below,font=\footnotesize] {$(2,0)$};
  \node[secondcolor] at (4,3.5) {$C_{f'}$};
\end{tikzpicture}
\end{center}

\begin{erwthma}
  \item Να μελετήσετε την $f$ ως προς τη μονοτονία και τα τοπικά ακρότατα.
  \item Αν γνωρίζετε ότι $f(2) = 3$, να υπολογίσετε το όριο $\lim\limits_{x \to 2} \dfrac{f(x) - 3}{x^2 - 4}$.
  \item Να μελετήσετε την $f$ ως προς την κυρτότητα. Ποια είναι η τετμημένη του σημείου καμπής της $C_f$;
  \item Αν επιπλέον $f(-1) = 5$, να αποδείξετε ότι η εξίσωση $f(x) = 6$ είναι αδύνατη στο διάστημα $[-3, 2]$.
\end{erwthma}
\end{thema}
%=========== ΤΕΛΟΣ - 107ο ΘΕΜΑ Β ===========

%---------------------------------------------------------------

%=========== 108ο ΘΕΜΑ Β ===========
\begin{thema}{Β}
Στο σχήμα δίνεται η γραφική παράσταση της $f'$ στο $[0, 6]$.
\begin{center}
\begin{tikzpicture}[scale=0.7]
  \draw[help lines, gray!30] (-1,-1) grid (7,2);
  \draw[-latex] (-1.3,0) -- (7.3,0) node[right] {$x$};
  \draw[-latex] (0,-1.3) -- (0,2.3) node[above] {$y$};
  \foreach \x in {1,2,3,4,5,6}
    \draw (\x,0.1) -- (\x,-0.1) node[below,font=\footnotesize] {$\x$};
  \foreach \y in {1,2}
    \draw (0.1,\y) -- (-0.1,\y) node[left,font=\footnotesize] {$\y$};
  \draw[very thick, secondcolor, domain=0:6, samples=60] plot (\x, {0.12*(\x-3)*(\x-3)});
  \filldraw[secondcolor] (3,0) circle (2pt) node[above,font=\footnotesize] {$(3,0)$};
  \filldraw[secondcolor] (0,1.08) circle (2pt);
  \filldraw[secondcolor] (6,1.08) circle (2pt);
  \node[secondcolor] at (5.5,1.5) {$C_{f'}$};
\end{tikzpicture}
\end{center}

\begin{erwthma}
  \item Να εξετάσετε αν η $f$ παρουσιάζει τοπικό ακρότατο στο $x = 3$, αιτιολογώντας πλήρως την απάντησή σας.
  \item Να εξετάσετε αν η $f$ είναι γνησίως αύξουσα στο $[0, 6]$.
  \item Αν $f(3) = 2$, να υπολογίσετε το όριο $\lim\limits_{x \to 3} \dfrac{f(x) - 2}{x - 3}$ και στη συνέχεια το $\lim\limits_{x \to 3} \dfrac{f(x) - 2}{(x-3)^2}$.
  \item Αν $f(0) = 1$, να αποδείξετε ότι $f(x) \geq 1$ για κάθε $x \in [0, 6]$.
\end{erwthma}
\end{thema}
%=========== ΤΕΛΟΣ - 108ο ΘΕΜΑ Β ===========

%---------------------------------------------------------------

%=========== 109ο ΘΕΜΑ Β ===========
\begin{thema}{Β}
Στο παρακάτω σχήμα δίνεται η γραφική παράσταση $C_f$ μιας συνεχούς συνάρτησης $f:\mathbb{R}\to\mathbb{R}$, η οποία έχει οριζόντια ασύμπτωτη στο $-\infty$ την ευθεία $y=2$, και παρουσιάζει \textbf{μοναδικό ολικό ελάχιστο} στο σημείο $A(0, -1)$. Η $C_f$ τέμνει τον άξονα $x'x$ ακριβώς στα σημεία με τετμημένες $x=-2$ και $x=2$.

\begin{center}
\begin{tikzpicture}[scale=0.8]
  \draw[help lines, gray!30] (-5,-2) grid (6,4);
  \draw[-latex] (-5.3,0) -- (6.3,0) node[right] {$x$};
  \draw[-latex] (0,-2.3) -- (0,4.3) node[above] {$y$};
  \foreach \x in {-4,-3,-2,-1,1,2,3,4,5}
    \draw (\x,0.1) -- (\x,-0.1) node[below,font=\footnotesize] {$\x$};
  \foreach \y in {-1,1,2,3}
    \draw (0.1,\y) -- (-0.1,\y) node[left,font=\footnotesize] {$\y$};
  
  \draw[dashed, secondcolor, thick] (-5,2) -- (0,2) node[above right,font=\footnotesize] {$y=2$};
  \draw[dashed, secondcolor, thick] (1.5,-0.5) -- (6,4) node[right,font=\footnotesize] {$y=x-2$};
  
  \draw[very thick, blue!70!black, smooth, tension=0.5] 
    plot coordinates {(-5,1.9) (-4,1.75) (-3,1) (-2,0) (-1,-0.8) (0,-1) (1,-0.8) (2,0) (4,1.8) (6,3.9)};
    
  \filldraw[blue!70!black] (0,-1) circle (2pt) node[below right,font=\footnotesize] {$A(0,-1)$};
  \filldraw[blue!70!black] (-2,0) circle (2pt);
  \filldraw[blue!70!black] (2,0) circle (2pt);
\end{tikzpicture}
\end{center}

\begin{erwthma}
  \item Να υπολογίσετε τα παρακάτω όρια, δικαιολογώντας την απάντησή σας:
  \[ L_1 = \lim_{x \to -\infty} e^{f(x)} \qquad \text{και} \qquad L_2 = \lim_{x \to 0} \frac{1}{f(x) + 1} \]
  \item Αξιοποιώντας το γράφημα, να λύσετε την ανίσωση: $f\left( f(x) + 2 \right) < 0$.
  \item Αν η $C_f$ δέχεται πλάγια ασύμπτωτη στο $+\infty$ την ευθεία $y = x - 2$, να υπολογίσετε το όριο:
  \[ \lim_{x \to +\infty} \frac{x f(x) - x^2 + 5x}{2x + \hm{x}} \]
  \item Αν η $f$ είναι παραγωγίσιμη στο $(0, 2)$, να αποδείξετε ότι υπάρχει τουλάχιστον ένα $\xi \in (0, 2)$ στο οποίο η εφαπτομένη της $C_f$ σχηματίζει γωνία $\omega \in \left(0, \frac{\pi}{2}\right)$ με τον άξονα $x'x$, και μάλιστα $\ef{\omega} = \frac{1}{2}$.
\end{erwthma}
\end{thema}
%=========== ΤΕΛΟΣ - 109ο ΘΕΜΑ Β ===========

%---------------------------------------------------------------

%=========== 110ο ΘΕΜΑ Β ===========
\begin{thema}{Β}
Στο σχήμα δίνεται η $C_{f'}$ μιας συνάρτησης $f:[0, 5]\to\mathbb{R}$.

\begin{center}
\begin{tikzpicture}[scale=0.75]
  \draw[help lines, gray!30] (-1,-3) grid (6,4);
  \draw[-latex] (-1.3,0) -- (6.3,0) node[right] {$x$};
  \draw[-latex] (0,-3.3) -- (0,4.3) node[above] {$y$};
  \foreach \x in {1,2,3,4,5}
    \draw (\x,0.1) -- (\x,-0.1) node[below,font=\footnotesize] {$\x$};
  \foreach \y in {-2,-1,1,2,3}
    \draw (0.1,\y) -- (-0.1,\y) node[left,font=\footnotesize] {$\y$};
  \draw[very thick, secondcolor, smooth, tension=0.5]
    plot coordinates {(0,2) (1,2.5) (2,0) (3,-2) (4,0) (4.5,-0.5) (5,-1.5)};
  \filldraw[secondcolor] (2,0) circle (2pt) node[above right,font=\footnotesize] {$(2,0)$};
  \filldraw[secondcolor] (4,0) circle (2pt) node[above right,font=\footnotesize] {$(4,0)$};
  \node[secondcolor] at (5,2) {$C_{f'}$};
\end{tikzpicture}
\end{center}

\begin{erwthma}
  \item Μελετήστε την $f$ ως προς τη μονοτονία και τα ακρότατα.
  \item Πόσα και ποια σημεία καμπής παρουσιάζει η $f$?
  \item Αν $f(0) = 3$, μπορείτε να πείτε αν $f(2) > 3$ ή $f(2) < 3$;
  \item Να βρεθεί η εξίσωση εφαπτομένης στο $x_0 = 4$, δεδομένου ότι $f(4) = -1$.
\end{erwthma}
\end{thema}
%=========== ΤΕΛΟΣ - 110ο ΘΕΜΑ Β ===========

%---------------------------------------------------------------
% Ασκήσεις 11-13: Σύνθεση f∘g
%---------------------------------------------------------------

%=========== 111ο ΘΕΜΑ Β ===========
\begin{thema}{Β}
Στο παρακάτω σχήμα δίνονται οι γραφικές παραστάσεις δύο συναρτήσεων $f$ και $g$ ορισμένων στο $[-2, 4]$.
\begin{center}
\begin{tikzpicture}[scale=0.7]
  \draw[help lines, gray!30] (-3,-2) grid (5,5);
  \draw[-latex] (-3.3,0) -- (5.3,0) node[right] {$x$};
  \draw[-latex] (0,-2.3) -- (0,5.3) node[above] {$y$};
  \foreach \x in {-2,-1,1,2,3,4}
    \draw (\x,0.1) -- (\x,-0.1) node[below,font=\footnotesize] {$\x$};
  \foreach \y in {-1,1,2,3,4}
    \draw (0.1,\y) -- (-0.1,\y) node[left,font=\footnotesize] {$\y$};
  % f: piecewise linear
  \draw[very thick, maincolor!70!black] (-2,1) -- (0,3) -- (2,1) -- (4,4);
  \node[maincolor!70!black] at (-1.5,2.5) {$C_f$};
  % g: piecewise linear
  \draw[very thick, orange!50!black] (-2,0) -- (0,2) -- (2,0) -- (4,2);
  \node[orange!50!black] at (3.5,0.5) {$C_g$};
  % key points
  \filldraw[maincolor!70!black] (0,3) circle (2pt) node[above,font=\footnotesize] {$(0,3)$};
  \filldraw[maincolor!70!black] (2,1) circle (2pt);
  \filldraw[orange!50!black] (0,2) circle (2pt) node[right,font=\footnotesize] {$(0,2)$};
  \filldraw[orange!50!black] (2,0) circle (2pt);
\end{tikzpicture}
\end{center}

\begin{erwthma}
  \item Να υπολογίσετε τα $(f \circ g)(0)$, $(f \circ g)(2)$ και $(g \circ f)(0)$.
  \item Πότε ορίζεται η $(f\circ g)(x)$; Να βρείτε το πεδίο ορισμού της.
  \item Να υπολογίσετε $(f \circ g)'(0)$, αν γνωρίζετε τις κλίσεις $f$ και $g$ στα κατάλληλα σημεία.
  \item Αν $h(x) = f(x) + g(x)$, να υπολογίσετε $h(0)$, $h(2)$ και να σχεδιάσετε πρόχειρα τη $C_h$ στο $[0, 2]$.
\end{erwthma}
\end{thema}
%=========== ΤΕΛΟΣ - 111ο ΘΕΜΑ Β ===========

%---------------------------------------------------------------

%=========== 112ο ΘΕΜΑ Β ===========
\begin{thema}{Β}
Δίνονται οι γραφικές παραστάσεις των συνεχών συναρτήσεων $f$ και $g$ στο $[0, 5]$.

\begin{center}
\begin{tikzpicture}[scale=0.7]
  \draw[help lines, gray!30] (-1,-1) grid (6,5);
  \draw[-latex] (-1.3,0) -- (6.3,0) node[right] {$x$};
  \draw[-latex] (0,-1.3) -- (0,5.3) node[above] {$y$};
  \foreach \x in {1,2,3,4,5}
    \draw (\x,0.1) -- (\x,-0.1) node[below,font=\footnotesize] {$\x$};
  \foreach \y in {1,2,3,4}
    \draw (0.1,\y) -- (-0.1,\y) node[left,font=\footnotesize] {$\y$};
  % f
  \draw[very thick, maincolor!70!black] (0,1) -- (2,3) -- (4,1) -- (5,2);
  \node[maincolor!70!black] at (1,3) {$C_f$};
  % g
  \draw[very thick, secondcolor] (0,0) -- (1,2) -- (2,4) -- (3,3) -- (5,1);
  \node[secondcolor] at (4.5,2) {$C_g$};
  % key points
  \filldraw[maincolor!70!black] (0,1) circle (2pt);
  \filldraw[maincolor!70!black] (2,3) circle (2pt);
  \filldraw[secondcolor] (1,2) circle (2pt);
  \filldraw[secondcolor] (2,4) circle (2pt);
\end{tikzpicture}
\end{center}

\begin{erwthma}
  \item Να υπολογίσετε $(f \circ g)(1)$, $(g \circ f)(2)$ και $(f \circ f)(0)$.
  \item Να βρείτε $(f+g)(2)$, $(f \cdot g)(0)$ και $\frac{f}{g}(3)$.
  \item Ορίζεται η συνάρτηση $(g \circ f)(x)$ στο $x = 4$? Αν ναι, υπολογίστε την τιμή $g(f(4))$.
  \item Λύστε γραφικά την εξίσωση $f(x) = g(x)$.
\end{erwthma}
\end{thema}
%=========== ΤΕΛΟΣ - 112ο ΘΕΜΑ Β ===========

%---------------------------------------------------------------

%=========== 113ο ΘΕΜΑ Β ===========
\begin{thema}{Β}
Δίνονται σε πίνακα τιμές δύο συναρτήσεων $f,g$ ορισμένων στο $\{0,1,2,3,4\}$:

\begin{center}
\begin{myhtblr}{}
$x$ & $0$ & $1$ & $2$ & $3$ & $4$ \\
$f(x)$ & $3$ & $1$ & $4$ & $0$ & $2$ \\
$g(x)$ & $2$ & $0$ & $1$ & $4$ & $3$ \\
\end{myhtblr}
\end{center}

\begin{erwthma}
  \item Να υπολογίσετε τα $(f \circ g)(0)$, $(g \circ f)(1)$, $(f \circ f)(2)$, $(g \circ g)(3)$.
  \item Εξετάστε αν η συνάρτηση $f$ είναι $1-1$.
  \item Να βρείτε τον αντίστροφο πίνακα τιμών $f^{-1}$ (αν η $f$ αντιστρέφεται) και τα $f^{-1}(0)$, $f^{-1}(3)$.
  \item Να λύσετε την εξίσωση $f(g(x)) = 2$.
\end{erwthma}
\end{thema}
%=========== ΤΕΛΟΣ - 113ο ΘΕΜΑ Β ===========

%---------------------------------------------------------------
% Ασκήσεις 14-16: Αντίστροφη f^{-1}
%---------------------------------------------------------------

%=========== 114ο ΘΕΜΑ Β ===========
\begin{thema}{Β}
Στο σχήμα δίνεται η γραφική παράσταση $C_f$ μιας γνησίως αύξουσας συνάρτησης $f:[0, 4]\to\mathbb{R}$ μαζί με τη {διχοτόμο} $y = x$.

\begin{center}
\begin{tikzpicture}[scale=0.7]
  \draw[help lines, gray!30] (-3,-3) grid (5,5);
  \draw[-latex] (-3.3,0) -- (5.3,0) node[right] {$x$};
  \draw[-latex] (0,-3.3) -- (0,5.3) node[above] {$y$};
  \foreach \x in {-2,-1,1,2,3,4}
    \draw (\x,0.1) -- (\x,-0.1) node[below,font=\footnotesize] {$\x$};
  \foreach \y in {-2,-1,1,2,3,4}
    \draw (0.1,\y) -- (-0.1,\y) node[left,font=\footnotesize] {$\y$};
  % y=x
  \draw[dashed, gray!60] (-3,-3) -- (5,5);
  \node[gray!60,font=\footnotesize] at (4.5,4.5) {$y=x$};
  % f
  \draw[very thick, maincolor!70!black] (0,-2) -- (1,0) -- (2,1) -- (3,2) -- (4,3);
  \node[maincolor!70!black] at (0.5,1.5) {$C_f$};
  \filldraw[maincolor!70!black] (0,-2) circle (2pt);
  \filldraw[maincolor!70!black] (1,0) circle (2pt);
  \filldraw[maincolor!70!black] (2,1) circle (2pt);
  \filldraw[maincolor!70!black] (3,2) circle (2pt);
  \filldraw[maincolor!70!black] (4,3) circle (2pt);
\end{tikzpicture}
\end{center}

\begin{erwthma}
  \item Εξηγήστε γιατί η $f$ είναι αντιστρέψιμη και να βρείτε το πεδίο ορισμού και το σύνολο τιμών της $f^{-1}$.
  \item Να βρείτε τα $f^{-1}(0)$, $f^{-1}(1)$, $f^{-1}(3)$.
  \item Αν $\dint_{0}^{4} f(x)\,dx = 2$, να υπολογίσετε το $\dint_{-2}^{3} f^{-1}(x)\,dx$.
  \item Ποια είναι η κλίση της εφαπτομένης της $C_{f^{-1}}$ στο σημείο με τετμημένη $x = 1$? Αιτιολογήστε τη απάντησή σας.
\end{erwthma}
\end{thema}
%=========== ΤΕΛΟΣ - 114ο ΘΕΜΑ Β ===========

%---------------------------------------------------------------

%=========== 115ο ΘΕΜΑ Β ===========
\begin{thema}{Β}
Δίνεται η γραφική παράσταση $C_f$ μιας \textbf{γνησίως φθίνουσας} συνάρτησης $f:[0, 4]\to[-1, 3]$.
\begin{center}
\begin{tikzpicture}[scale=0.8]
  \draw[help lines, gray!30] (-2,-2) grid (5,4);
  \draw[-latex] (-2.3,0) -- (5.3,0) node[right] {$x$};
  \draw[-latex] (0,-2.3) -- (0,4.3) node[above] {$y$};
  \foreach \x in {-1,1,2,3,4}
    \draw (\x,0.1) -- (\x,-0.1) node[below,font=\footnotesize] {$\x$};
  \foreach \y in {-1,1,2,3}
    \draw (0.1,\y) -- (-0.1,\y) node[left,font=\footnotesize] {$\y$};
  \draw[dashed, gray!60] (-2,-2) -- (4,4);
  \draw[very thick, maincolor!70!black] (0,3) -- (1,2) -- (2,1) -- (3,0) -- (4,-1);
  \node[maincolor!70!black] at (1,3) {$C_f$};
  \filldraw[maincolor!70!black] (0,3) circle (2pt);
  \filldraw[maincolor!70!black] (3,0) circle (2pt);
  \filldraw[maincolor!70!black] (4,-1) circle (2pt);
\end{tikzpicture}
\end{center}

\begin{erwthma}
  \item Δείξτε ότι η $f$ αντιστρέφεται και βρείτε το πεδίο ορισμού της $f^{-1}$;
  \item Να βρείτε τα $f^{-1}(-1)$, $f^{-1}(0)$, $f^{-1}(2)$.
  \item Η $f^{-1}$ είναι γνησίως αύξουσα ή γνησίως φθίνουσα?
  \item Παρατηρώντας το σχήμα, να δείξετε ότι $f^{-1}(x) = f(x)$ για κάθε $x \in [-1, 3]$, δηλαδή ότι η $f$ ταυτίζεται με την αντίστροφή της.
  \item Αν $\dint_{0}^{4} f(x)\,dx = 4$, να υπολογίσετε το $\dint_{-1}^{3} f^{-1}(x)\,dx$.
\end{erwthma}
\end{thema}
%=========== ΤΕΛΟΣ - 115ο ΘΕΜΑ Β ===========

%---------------------------------------------------------------

%=========== 116ο ΘΕΜΑ Β ===========
\begin{thema}{Β}
Στο παρακάτω σχήμα δίνονται οι γραφικές παραστάσεις $C_f$ {και} η $C_{f^{-1}}$ των συναρτήσεων $f$ και $f^{-1}$ στο διάστημα $[0,5]$.

\begin{center}
\begin{tikzpicture}[scale=0.7]
  \draw[help lines, gray!30] (-1,-1) grid (6,6);
  \draw[-latex] (-1.3,0) -- (6.3,0) node[right] {$x$};
  \draw[-latex] (0,-1.3) -- (0,6.3) node[above] {$y$};
  \foreach \x in {1,2,3,4,5}
    \draw (\x,0.1) -- (\x,-0.1) node[below,font=\footnotesize] {$\x$};
  \foreach \y in {1,2,3,4,5}
    \draw (0.1,\y) -- (-0.1,\y) node[left,font=\footnotesize] {$\y$};
  \draw[dashed, gray!60] (-1,-1) -- (6,6);
  % f: concave curve below y=x
  \draw[very thick, maincolor!70!black, smooth, tension=0.5]
    plot coordinates {(0,0) (1,0.5) (2,1.2) (3,2.2) (4,3.5) (5,5)};
  \node[maincolor!70!black] at (1.5,2.5) {$C_f$};
  % f^{-1}: convex curve above y=x (mirror)
  \draw[very thick, secondcolor, smooth, tension=0.5]
    plot coordinates {(0,0) (0.5,1) (1.2,2) (2.2,3) (3.5,4) (5,5)};
  \node[secondcolor] at (4,2.5) {$C_{f^{-1}}$};
  \filldraw[maincolor!70!black] (0,0) circle (2pt);
  \filldraw[maincolor!70!black] (5,5) circle (2pt) node[above left,font=\footnotesize] {$(5,5)$};
  \filldraw[maincolor!70!black] (2,1.2) circle (2pt) node[below right,font=\footnotesize] {$(2, 1{,}2)$};
  \filldraw[secondcolor] (1.2,2) circle (2pt) node[above left,font=\footnotesize] {$(1{,}2, 2)$};
\end{tikzpicture}
\end{center}

\begin{erwthma}
  \item Αν $f'(2) = \dfrac{3}{5}$, να υπολογίσετε $(f^{-1})'(1{,}2)$ και να γράψετε την εξίσωση της εφαπτομένης της $C_{f^{-1}}$ στο σημείο $A'(1{,}2,\; 2)$.
  \item Παρατηρώντας το σχήμα, να αποδείξετε ότι $f(x) \leq x$ για κάθε $x \in [0, 5]$, με ισότητα μόνο στα άκρα. Στη συνέχεια, να συμπεράνετε την αντίστοιχη ανισότητα για την $f^{-1}$.
  \item Χρησιμοποιώντας αποκλειστικά το σχήμα, να αποδείξετε ότι:
  \[ \dint_0^5 f(x)\,dx + \dint_0^5 f^{-1}(x)\,dx = 25. \]
  \item Αν $E$ είναι το εμβαδόν του χωρίου που περικλείεται μεταξύ $C_f$ και $C_{f^{-1}}$, να εκφράσετε το $\dint_0^5 f(x)\,dx$ συναρτήσει του $E$.
\end{erwthma}
\end{thema}
%=========== ΤΕΛΟΣ - 116ο ΘΕΜΑ Β ===========

%---------------------------------------------------------------
% Ασκήσεις 17-20: Μικτές (δυσκολότερες Β)
%---------------------------------------------------------------

%=========== 117ο ΘΕΜΑ Β ===========
\begin{thema}{Β}
Δίνεται η γραφική παράσταση $C_f$ μιας συνεχούς $f:[-3, 3]\to\mathbb{R}$ η οποία είναι \textbf{άρτια} ($f(-x) = f(x)$). Στο σχήμα φαίνεται μόνο η $C_f$ για $x \geq 0$, με $f(0) = 4$, $f(1) = 2$, $f(2) = 0$, $f(3) = 1$.

\begin{center}
\begin{tikzpicture}[scale=0.7]
  \draw[help lines, gray!30] (-4,-1) grid (4,5);
  \draw[-latex] (-4.3,0) -- (4.3,0) node[right] {$x$};
  \draw[-latex] (0,-1.3) -- (0,5.3) node[above] {$y$};
  \foreach \x in {-3,-2,-1,1,2,3}
    \draw (\x,0.1) -- (\x,-0.1) node[below,font=\footnotesize] {$\x$};
  \foreach \y in {1,2,3,4}
    \draw (0.1,\y) -- (-0.1,\y) node[left,font=\footnotesize] {$\y$};
  % right half
  \draw[very thick, maincolor!70!black] (0,4) -- (1,2) -- (2,0) -- (3,1);
  \filldraw[maincolor!70!black] (0,4) circle (2pt) node[above right,font=\footnotesize] {$(0,4)$};
  \filldraw[maincolor!70!black] (2,0) circle (2pt);
  \filldraw[maincolor!70!black] (-2,0) circle (2pt);
  \node[maincolor!70!black,font=\footnotesize] at (2.5,2) {$C_f$};
\end{tikzpicture}
\end{center}

\begin{erwthma}
  \item Αξιοποιώντας το γράφημα για να διαβάσετε την κλίση της $C_f$ στο $x = 2$, να υπολογίσετε το όριο:
  \[ \lim\limits_{x \to 2} \frac{f(x)}{x - 2}. \]
  \item Να αποδείξετε ότι η εξίσωση $f(x) = 1$ έχει τουλάχιστον μία ρίζα στο διάστημα $(2, 3)$, εφαρμόζοντας το Θεώρημα \eng{Bolzano} σε κατάλληλη συνάρτηση.
  \item Η $f$ παραγωγίσιμη στο $[-2, 2]$. Να αποδείξετε ότι υπάρχει $\xi \in (-2, 2)$ με $f'(\xi) = 0$, εφαρμόζοντας το Θεώρημα \eng{Rolle}. Ποια είναι η γεωμετρική ερμηνεία του $\xi$ στο σχήμα?
  \item Αν η $f$ είναι παραγωγίσιμη και $f'(x) < 0$ για κάθε $x \in (0, 2)$, να αποδείξετε ότι $f$ δεν έχει ρίζες στο $(0, 2)$, χρησιμοποιώντας τη μονοτονία και την τιμή $f(0) = 4$.
\end{erwthma}
\end{thema}
%=========== ΤΕΛΟΣ - 117ο ΘΕΜΑ Β ===========

%---------------------------------------------------------------

%=========== 118ο ΘΕΜΑ Β ===========
\begin{thema}{Β}
Δίνεται η γραφική παράσταση $C_f$ μιας γνησίως αύξουσας $f:[0, 4]\to\mathbb{R}$ με $f(0) = 0$. Δίνεται επιπλέον ότι η σκιασμένη περιοχή $E_1$ (μεταξύ $C_f$ και άξονα $y$ από $y=0$ ως $y=f(4)=3$) έχει εμβαδόν $E_1 = 5$.

\begin{center}
\begin{tikzpicture}[scale=0.85]
  \draw[help lines, gray!30] (-1,-1) grid (5,4);
  \draw[-latex] (-1.3,0) -- (5.3,0) node[right] {$x$};
  \draw[-latex] (0,-1.3) -- (0,4.3) node[above] {$y$};
  \foreach \x in {1,2,3,4}
    \draw (\x,0.1) -- (\x,-0.1) node[below,font=\footnotesize] {$\x$};
  \foreach \y in {1,2,3}
    \draw (0.1,\y) -- (-0.1,\y) node[left,font=\footnotesize] {$\y$};
  % f
  \draw[very thick, maincolor!70!black, smooth, tension=0.5]
    plot coordinates {(0,0) (1,0.5) (2,1.2) (3,2.2) (4,3)};
  % shaded area under f
  \fill[cyan!15, opacity=0.5, smooth, tension=0.5]
    (0,0) -- plot coordinates {(0,0) (1,0.5) (2,1.2) (3,2.2) (4,3)} -- (4,0) -- cycle;
  \node at (2,0.5) {$E_2$};
  % shaded area between f and y=3 (i.e. area to the left of f)
  \fill[yellow!20, opacity=0.5, smooth, tension=0.5]
    (0,0) -- plot coordinates {(0,0) (1,0.5) (2,1.2) (3,2.2) (4,3)} -- (4,3) -- (0,3) -- (0,0);
  \node at (1,2.5) {$E_1 = 5$};
  \filldraw[maincolor!70!black] (0,0) circle (2pt);
  \filldraw[maincolor!70!black] (4,3) circle (2pt) node[above right,font=\footnotesize] {$(4,3)$};
  \node[maincolor!70!black] at (3,0.8) {$C_f$};
  \draw[dashed] (4,0) -- (4,3) -- (0,3);
\end{tikzpicture}
\end{center}

\begin{erwthma}
  \item Να υπολογίσετε το $\dint_0^4 f(x)\,dx$, χρησιμοποιώντας τη γεωμετρική σχέση $E_1 + E_2 = 4 \cdot f(4)$ που φαίνεται στο σχήμα. Τι αντιπροσωπεύει το γινόμενο $4 \cdot f(4)$?
  \item Να αποδείξετε ότι η εξίσωση $f(x) = 2$ έχει τουλάχιστον μία λύση στο ανοικτό διάστημα $(0, 4)$, εφαρμόζοντας το Θεώρημα \eng{Bolzano} σε κατάλληλη συνάρτηση.
  \item Αν η $f$ είναι παραγωγίσιμη, να αποδείξετε ότι υπάρχει $\xi \in (0, 4)$ τέτοιο ώστε $f'(\xi) = \dfrac{3}{4}$. Ποια η γεωμετρική ερμηνεία αυτού του $\xi$ στο σχήμα?
  \item Αν $g(x) = \dint_0^{f(x)} f^{-1}(t)\,dt$, να υπολογίσετε το $g(4)$ χρησιμοποιώντας τα δεδομένα του σχήματος.
\end{erwthma}
\end{thema}
%=========== ΤΕΛΟΣ - 118ο ΘΕΜΑ Β ===========

%---------------------------------------------------------------

%=========== 119ο ΘΕΜΑ Β ===========
\begin{thema}{Β}
Δίνεται η γραφική παράσταση $C_f$ μιας συνεχούς $f:[-2, 3]\to\mathbb{R}$, μαζί μe τις εφαπτόμενές της σε δύο σημεία $A(0, 2)$ και $B(2, 0)$. Η εφαπτομένη στο $A$ έχει κλίση $1$ και στο $B$ κλίση $-2$.

\begin{center}
\begin{tikzpicture}[scale=0.8]
  \draw[help lines, gray!30] (-3,-2) grid (4,4);
  \draw[-latex] (-3.3,0) -- (4.3,0) node[right] {$x$};
  \draw[-latex] (0,-2.3) -- (0,4.3) node[above] {$y$};
  \foreach \x in {-2,-1,1,2,3}
    \draw (\x,0.1) -- (\x,-0.1) node[below,font=\footnotesize] {$\x$};
  \foreach \y in {-1,1,2,3}
    \draw (0.1,\y) -- (-0.1,\y) node[left,font=\footnotesize] {$\y$};
  % curve
  \draw[very thick, maincolor!70!black, smooth, tension=0.5]
    plot coordinates {(-2,0) (-1,0.8) (0,2) (1,2.5) (2,0) (3,-1)};
  % tangent at A(0,2), slope 1
  \draw[dashed, secondcolor] (-1,1) -- (1.5,3.5);
  \node[secondcolor,font=\footnotesize] at (1.5,3.8) {$\varepsilon_A$};
  % tangent at B(2,0), slope -2
  \draw[dashed, secondcolor] (1,2) -- (3,-2);
  \node[secondcolor,font=\footnotesize] at (3,-1.7) {$\varepsilon_B$};
  \filldraw[maincolor!70!black] (0,2) circle (2pt) node[above left,font=\footnotesize] {$A(0,2)$};
  \filldraw[maincolor!70!black] (2,0) circle (2pt) node[below left,font=\footnotesize] {$B(2,0)$};
\end{tikzpicture}
\end{center}

\begin{erwthma}
  \item Να εφαρμόσετε το Θεώρημα \eng{Rolle} στο $[-2, 2]$ και να αποδείξετε ότι υπάρχει $\xi \in (-2, 2)$ με $f'(\xi) = 0$. Ποια η γεωμετρική ερμηνεία του $\xi$ στο σχήμα?
  \item Γράψτε τις εξισώσεις των εφαπτομένων $\varepsilon_A$ και $\varepsilon_B$.
  \item Αξιοποιώντας τα δεδομένα του σχήματος, να υπολογίσετε τα όρια:
  \[ \lim\limits_{x \to 0} \frac{f(x) - 2}{x} \qquad \text{και} \qquad \lim\limits_{x \to 2} \frac{f(x)}{x - 2}. \]
  \item Αν $h(x) = f(x) \cdot x$, να βρείτε $h'(0)$ και $h'(2)$.
\end{erwthma}
\end{thema}
%=========== ΤΕΛΟΣ - 119ο ΘΕΜΑ Β ===========

%---------------------------------------------------------------

%=========== 120ο ΘΕΜΑ Β ===========
\begin{thema}{Β}
Στο σχήμα δίνεται η γραφική παράσταση $C_{f''}$ της \textbf{δεύτερης παραγώγου} μιας δις παραγωγίσιμης συνάρτησης $f:[-3, 3]\to\mathbb{R}$.

\begin{center}
\begin{tikzpicture}[scale=0.8]
  \draw[help lines, gray!30] (-4,-3) grid (4,3);
  \draw[-latex] (-4.3,0) -- (4.3,0) node[right] {$x$};
  \draw[-latex] (0,-3.3) -- (0,3.3) node[above] {$y$};
  \foreach \x in {-3,-2,-1,1,2,3}
    \draw (\x,0.1) -- (\x,-0.1) node[below,font=\footnotesize] {$\x$};
  \foreach \y in {-2,-1,1,2}
    \draw (0.1,\y) -- (-0.1,\y) node[left,font=\footnotesize] {$\y$};
  
  \draw[very thick, maincolor!70!black, smooth, domain=-3.1:3.1] plot (\x, {-0.5*(\x*\x - 4)});
  \node[maincolor!70!black] at (1.5, 1.8) {$C_{f''}$};
  \filldraw[maincolor!70!black] (0,2) circle (2pt) node[above right,font=\footnotesize] {$(0,2)$};
  \filldraw[maincolor!70!black] (-2,0) circle (2pt);
  \filldraw[maincolor!70!black] (2,0) circle (2pt);
\end{tikzpicture}
\end{center}

\begin{erwthma}
  \item Να μελετήσετε την $f$ ως προς την κυρτότητα και να βρείτε τις τετμημένες των σημείων καμπής της $C_f$.
  \item Δίνεται ότι το εμβαδόν του χωρίου που περικλείεται από την $C_{f''}$ και τον άξονα $x'x$ είναι ίσο με $\dfrac{16}{3}$. Αν η εφαπτομένη της $C_f$ στο $x=-2$ είναι παράλληλη στην ευθεία $y = -x + 2024$, να βρείτε την κλίση της εφαπτομένης της $C_f$ στο $x=2$.
  \item Αν γνωρίζετε επιπλέον ότι $f'(0) = 0$, να μελετήσετε την $f$ ως προς τη μονοτονία στο διάστημα $[-2, 2]$ και να προσδιορίσετε το είδος του τοπικού ακροτάτου στο $x=0$.
  \item Με δεδομένο ότι $f'(0)=0$, να υπολογίσετε το όριο:
  \[ \lim\limits_{x \to 0} \frac{f'(x)}{x} \]
\end{erwthma}
\end{thema}
%=========== ΤΕΛΟΣ - 120ο ΘΕΜΑ Β ===========